\documentclass[11pt,a4paper]{article}

\usepackage[T1]{fontenc}
\usepackage[utf8]{inputenc}
\usepackage{lmodern}
\usepackage{microtype}
\usepackage[a4paper,margin=2.45cm]{geometry}
\usepackage{amsmath,amssymb,amsthm,mathtools}
\usepackage{bm}
\usepackage{booktabs}
\usepackage{array}
\usepackage{enumitem}
\usepackage{xcolor}
\usepackage{url}
\usepackage{hyperref}

\definecolor{linkblue}{RGB}{20,68,126}
\hypersetup{
  colorlinks=true,
  linkcolor=linkblue,
  citecolor=linkblue,
  urlcolor=linkblue,
  pdftitle={Gauge Freedom in Discrete Stokes Structures},
  pdfauthor={Redha Moulla},
  pdfsubject={Abstract discrete Stokes structures, pairing gauges, kernel polynomials, and Gauss--Legendre rigidity},
  pdfkeywords={discrete Stokes structure, Stokes--Dirac structure, pairing gauge, kernel polynomial, Gauss--Legendre quadrature, summation by parts, port-Hamiltonian system}
}

\setlist[itemize]{leftmargin=*,topsep=4pt,itemsep=2pt}
\setlist[enumerate]{leftmargin=*,topsep=4pt,itemsep=2pt}
\allowdisplaybreaks

\newtheorem{theorem}{Theorem}[section]
\newtheorem{proposition}[theorem]{Proposition}
\newtheorem{corollary}[theorem]{Corollary}
\newtheorem{lemma}[theorem]{Lemma}
\theoremstyle{definition}
\newtheorem{definition}[theorem]{Definition}
\newtheorem{example}[theorem]{Example}
\theoremstyle{remark}
\newtheorem{remark}[theorem]{Remark}

\newcommand{\R}{\mathbb{R}}
\newcommand{\Pp}{\mathbb{P}}
\newcommand{\Ker}{\operatorname{ker}}
\newcommand{\rank}{\operatorname{rank}}
\newcommand{\im}{\operatorname{im}}
\newcommand{\Rad}{\operatorname{rad}}
\newcommand{\spanof}{\operatorname{span}}
\newcommand{\diag}{\operatorname{diag}}
\newcommand{\cA}{\mathcal{A}}
\newcommand{\cB}{\mathcal{B}}
\newcommand{\cC}{\mathcal{C}}
\newcommand{\cD}{\mathcal{D}}
\newcommand{\cG}{\mathcal{G}}
\newcommand{\cK}{\mathcal{K}}
\newcommand{\cN}{\mathcal{N}}
\newcommand{\cP}{\mathcal{P}}
\newcommand{\dd}{\,\mathrm{d}}
\newcommand{\pair}[2]{\left\langle #1,#2\right\rangle}
\newcommand{\bracket}[1]{\left[#1\right]_{-1}^{1}}

\title{\bfseries Gauge Freedom in Discrete Stokes Structures:\\
Affine Pairings, Intrinsic Nodes, and the\\
Rigidity of Gauss--Legendre Locality}
\author{Redha Moulla\\
\small AXIA, France\\
\small \href{mailto:redha.moulla@axia-conseil.com}{redha.moulla@axia-conseil.com}}
\date{August 2026}

\begin{document}
\maketitle

\begin{abstract}
We define a discrete Stokes structure by a surjection $D:E\to F$ and a symmetric boundary
form $b$.  Compatible mixed pairings exist exactly when
$b|_{\Ker D\times\Ker D}=0$; when nonempty, they form an affine space directed by
$\Lambda^2(E/\Ker D)^*$.  We derive the full-rank criterion and classify left-radical
lines for $\dim\Ker D=1$ and $\rank b=2$.

For the polynomial complex $\Pp_N\xrightarrow{d/dx}\Pp_{N-1}$, we determine the exact image,
topology, and fibres of the projective kernel map.  Interior simple roots are admissible
exactly when $\sum_i\operatorname{artanh}(x_i)=0$.  Each full-rank fibre has dimension
$(N-1)(N-2)/2$ and, for $N\geq3$, two Pfaffian components.

Scalar localization forces exactness through degree $2N-1$ and uniquely selects the zero
gauge, Gauss--Legendre nodes, and positive Gauss weights.
Thus Gauss--Legendre is rigid under locality but not fixed by Stokes alone.

With the $L^2$ quotient metric, the canonical prescribed-kernel representative has rank $N$
in exact degree and $N-2$ otherwise.  The infimum over the full-rank fibre is attained if and
only if the kernel has degree $N$, and the lower-degree defect opens linearly under a Legendre
perturbation.  A global continuous, equivalently smooth, full-rank synthesis exists exactly
when $N$ is odd or $N\in\{2,4,8\}$; otherwise an explicit two-chart Rodrigues atlas is optimal.
In even degree, every full-rank synthesis diverges near the deleted constant class.  We
relate the gauge to rectangular SBP and quotient-compatible skew operators on $H^1/\R$,
while separating it from local $P^{(0)}$ terms, boundary-port changes, and weighted metrics.
\end{abstract}

\noindent\textbf{Keywords.}
discrete Stokes structure; Stokes--Dirac structure; mixed pairing; gauge freedom; kernel polynomial; Gauss--Legendre
quadrature; summation by parts; port-Hamiltonian system; orthogonal polynomials.

\noindent\textbf{MSC 2020.}
65D32; 65M12; 37J39; 46N20; 93C20.

\section{Introduction}

Stokes' theorem fixes a boundary power balance, but it does not by itself prescribe a metric,
a quadrature rule, or a nodal grid.  This distinction is fundamental in open
port-Hamiltonian systems.  The Stokes--Dirac structure specifies the power-conserving
interconnection and its boundary ports, whereas the Hamiltonian density or constitutive map
specifies how energy variables are converted into efforts
\cite{vanderschaftmaschke2002,legorrec2005}.

The geometric pseudospectral construction of Moulla, Lef\`evre, and Maschke
\cite{moulla2012} realizes this separation with two polynomial spaces of unequal dimensions:
efforts belong to $\Pp_N$, flows to $\Pp_{N-1}$, differentiation is exact between them, and
the power pairing is represented by a rectangular matrix.  The construction is valid in
arbitrary unisolvent flow coordinates.  For the exact unweighted $L^2$ pairing, its
one-dimensional radical is generated by the degree-$N$ Legendre polynomial, and, as proved
below, a same-grid diagonal power chart is possible precisely at the Gauss--Legendre nodes.

Power-preserving projection and dual-representation choices also occur in weak-form and
dual-field finite-element discretizations of Stokes--Dirac systems
\cite{kotyczka2018,brugnolirashad2022}.  Those works construct particular geometric
discretizations.  The question addressed here is complementary: with the length-one Stokes
data $(E,F,D,b)$ held fixed, classify the entire antisymmetric bilinear freedom and determine
how it moves the polynomial radical.

The present paper asks a different question.  Instead of fixing the $L^2$ pairing, we retain
only the exact Stokes identity and classify every mixed pairing compatible with it.  We first
solve this problem for an arbitrary finite-dimensional surjective complex $E\xrightarrow D F$
with a boundary form, and only then specialize to polynomials.  The difference between two
compatible pairings, after composition with $D$, is a two-form on $E/\Ker D$.  We call this freedom a
\emph{pairing gauge}.  The term is used in the precise linear-algebraic sense of a closed
two-form or $B$-field transform; it should not be confused with a change of polynomial basis,
with the gauge symmetry of physical potentials, or with a change of Hamiltonian density.

Once the gauge is exposed, the radical becomes gauge-dependent.  Its generator
$\kappa_\beta$ may have roots different from the Legendre points.  The central inverse
problem is therefore:
\begin{quote}
Which projective polynomials, and which real node configurations, can occur as the radical
of a full-rank Stokes-compatible mixed pairing?
\end{quote}
We solve this problem exactly.  The image lies in a projective hyperplane and is obtained by
removing explicit algebraic degeneracy loci: endpoint-vanishing lines, and, in even degree,
the constant line.  It is therefore a relative semialgebraic open subset of that hyperplane,
not the whole hyperplane.  On the degree-$N$ stratum, the corresponding root sets satisfy one
nonlinear but elementary balance law.  This produces a codimension-one family of admissible
grids, rather than all grids.

The second question is whether these roots are quadrature nodes in the classical sense.  The
answer is negative in general and rigid in the only positive case.  Every full-rank pairing
can be represented at its distinct kernel roots by an invertible matrix of nodal couplings.
That matrix is diagonal if and only if the gauge vanishes.  Stokes compatibility and scalar
locality together force exactness through degree $2N-1$, and hence force the ordinary
Gauss--Legendre rule.  This result isolates the precise content of Gauss--Legendre
canonicity: it is not absolute, but it is rigid once the unweighted Stokes geometry and local
scalar power evaluation are required.

The logical dependence of the result is
\begin{equation}\label{eq:logicaldiagram}
 \begin{gathered}
 \boxed{\text{Stokes data}}\Longrightarrow
 \boxed{\text{affine pairing class}}\Longrightarrow
 \boxed{\text{intrinsic radical}}\Longrightarrow
 \boxed{\text{full nodal matrix}},\\[1mm]
 \boxed{\text{scalar locality}}\Longrightarrow
 \boxed{L^2}\Longrightarrow
 \boxed{\text{Gauss--Legendre}}.
 \end{gathered}
\end{equation}
The first chain holds throughout the full-rank stratum; the second is the additional rigidity
statement.

The main contributions are as follows.
\begin{enumerate}
\item We prove the abstract affine-classification theorem, a necessary and sufficient
full-rank existence criterion, and a radical-line theorem for rank-two boundary forms.
\item We determine the exact polynomial kernel image, its topology, construct every
admissible line, and compute both the dimension and topology of every full-rank fibre.  We
also classify the orders admitting a global continuous or smooth full-rank synthesis over
the whole image, determine the minimal atlas size, give an explicit two-chart formula, and
prove the unavoidable radial blow-up near the deleted constant class.
\item We characterize the simple interior root sets by a hyperbolic barycentre condition and
separate intrinsic kernel nodes, full matrix localization, and scalar quadrature.
\item We formulate positive localization as a truncated moment problem, use the differentiated
product span to prove Gauss--Legendre rigidity, and record the independent Tchakaloff
existence mechanism for positive moment functionals.
\item Relative to the $L^2$ quotient metric, we derive the closed formula for the
minimum-norm solution of a prescribed-kernel constraint and prove its global rank
stratification: its rank is $N$ for degree-$N$ kernels and exactly $N-2$ on every nonempty
lower-degree stratum, where both rectangular radicals are computed explicitly.  We prove
that the full-rank minimum is attained exactly in the degree-$N$ stratum, that the lower
defect opens transversely in the Legendre direction with two-sided singular-value bounds,
and then establish local
stability near Gauss and show why no uniform global estimate can ignore node collision.
\item We give the precise rectangularization statement for SBP, an exact continuous operator
model with a coherent-extension criterion, and a careful comparison with Dirac two-form
gauges, $P^{(0)}$ freedom, boundary-port coordinates, and weighted metrics.
\end{enumerate}

The relation between SBP and quadrature is classical
\cite{hickenzingg2013,delrey2014}; function-space SBP extends the theory to non-polynomial
spaces \cite{glaubitz2023}, and recent work constructs minimal FSBP operators through
generalized Gaussian or Gauss--Lobatto rules \cite{hale2026,bercik2026}.  Gauge equivalence
of Dirac structures by closed two-forms is also established in Dirac geometry
\cite{bursztynradko2003}.  The contribution claimed here is narrower and different: an exact
classification of the abstract affine pairing freedom and, in the polynomial instance, of the
projective kernel map, together with the locality rigidity and local stability theorems.

\section{Abstract discrete Stokes structures}\label{sec:abstract}

\subsection{Definition and affine classification}

Let $E$ and $F$ be finite-dimensional real vector spaces, let
$D:E\to F$ be surjective, and let $b\in S^2E^*$ be a symmetric bilinear form.  Put
$K=\Ker D$ and $m=\dim F$.

\begin{definition}[Discrete Stokes structure and compatible pairing]
The quadruple $(E,F,D,b)$ is called a \emph{discrete Stokes structure}.  A mixed bilinear
pairing $M:E\times F\to\R$ is \emph{Stokes-compatible} if
\begin{equation}\label{eq:abstractstokes}
 M(e,De')+M(e',De)=b(e,e'),\qquad e,e'\in E.
\end{equation}
The set of all such pairings is denoted by $\cC(D,b)$.
A compatible pairing is \emph{full rank} when the map
$M^\flat:E\to F^*$, $e\mapsto M(e,\cdot)$, has rank $m=\dim F$.
\end{definition}

Associate with $M$ the bilinear form
\begin{equation}\label{eq:abstracta}
 a_M(e,e')=M(e,De').
\end{equation}
It satisfies $a_M+a_M^{\mathsf T}=b$ and $a_M(E,K)=0$.

\begin{theorem}[Abstract affine classification]\label{thm:abstractaffine}
The class $\cC(D,b)$ is nonempty if and only if
\begin{equation}\label{eq:abstractobstruction}
 b|_{K\times K}=0.
\end{equation}
When this condition holds, $\cC(D,b)$ is an affine space modelled on
\begin{equation}\label{eq:abstractgauge}
 \Lambda^2(E/K)^*\simeq\Lambda^2F^*,
\end{equation}
and therefore
\begin{equation}\label{eq:abstractdimension}
 \dim\cC(D,b)=\frac{m(m-1)}2.
\end{equation}
\end{theorem}

\begin{proof}
If $M$ is compatible and $k,k'\in K$, then both terms on the left of
\eqref{eq:abstractstokes} vanish, hence $b(k,k')=0$.

Conversely, choose a splitting $E=K\oplus H$.  Define $a:E\times E\to\R$ by
\[
 a(e,k)=0,\qquad a(k,h)=b(k,h),\qquad
 a(h,h')=\tfrac12b(h,h')
\]
for $e\in E$, $k\in K$, and $h,h'\in H$.  The assumption on $b$ gives
$a+a^{\mathsf T}=b$, and $a(E,K)=0$.  Therefore
\[
 M(e,f)=a(e,h_f),
\]
where $h_f\in H$ is the unique vector satisfying $Dh_f=f$, defines a compatible pairing.

If $M_1,M_0$ are compatible, then
$\omega=a_{M_1}-a_{M_0}$ is skew-symmetric and $\omega(E,K)=0$.  Skew-symmetry also gives
$\omega(K,E)=0$, so $\omega$ descends uniquely to $\Lambda^2(E/K)^*$.  Conversely, every
such two-form may be added to $a_{M_0}$.  Finally $E/K\simeq F$ because $D$ is surjective,
which proves \eqref{eq:abstractdimension}.
\end{proof}

The theorem separates an existence obstruction from the gauge freedom.  It also specifies
which radical statement is valid: $K$ is always contained in the right radical of $a_M$, but
\begin{equation}\label{eq:leftradicalK}
 a_M(k,e)=b(k,e),\qquad k\in K, e\in E,
\end{equation}
so $K$ need not lie in the left radical.

\begin{theorem}[Existence of full-rank compatible pairings]
\label{thm:abstractfullrank}
Assume $b|_{K\times K}=0$, so that $\cC(D,b)$ is nonempty.  A full-rank compatible pairing
exists if and only if
\begin{equation}\label{eq:abstractfullrankcriterion}
 \boxed{\ m\ \text{is even}\quad\text{or}\quad b\neq0.\ }
\end{equation}
Equivalently, the only obstruction beyond compatibility is the boundary-free odd-dimensional
case $b=0$ and $m$ odd.
\end{theorem}

\begin{proof}
Choose a splitting $E=K\oplus H$ and identify $H\simeq F$.  Relative to this splitting,
every compatible form $a_M$ has a matrix
\begin{equation}\label{eq:abstractblockform}
 [a_M]=
 \begin{pmatrix}
  0&B\\
  0&S+\Omega
 \end{pmatrix},
 \qquad S^{\mathsf T}=S,\quad \Omega^{\mathsf T}=-\Omega,
\end{equation}
where
$[b]=\left(\begin{smallmatrix}0&B\\B^{\mathsf T}&2S\end{smallmatrix}\right)$.
The ranks of $M$ and $a_M$ agree because $D$ is onto.  Hence $M$ is full rank precisely
when
\begin{equation}\label{eq:abstractinjectivity}
 H\longrightarrow K^*\oplus H^*,
 \qquad h\longmapsto\bigl(Bh,(S+\Omega)h\bigr)
\end{equation}
is injective.

If $m$ is even, choose a nondegenerate skew form $\Omega_0$ on $H$.  Then
$S+t\Omega_0$ is invertible for all sufficiently large $|t|$, so
\eqref{eq:abstractinjectivity} is injective.  Let now $m$ be odd and $b\neq0$.  If $B\neq0$,
choose $v\in H$ with $Bv\neq0$ and a skew form $\Omega_0$ with
$\Ker\Omega_0=\spanof\{v\}$.  The map obtained from $\Omega=t\Omega_0$ is injective for
all sufficiently large $|t|$: otherwise unit null vectors along a sequence
$|t|\to\infty$ would converge to a nonzero multiple of $v$ while still belonging to
$\Ker B$, a contradiction.  If $B=0$, then $S\neq0$.  Choose $v$ with
$S(v,v)\neq0$ and again $\Ker\Omega_0=\spanof\{v\}$.  A hypothetical sequence of unit
vectors in $\Ker(S+t\Omega_0)$ converges to a multiple of $v$; pairing the kernel equation
with $v$ and using $\Omega_0v=0$ contradicts $S(v,v)\neq0$.

Finally, if $b=0$, then $B=S=0$ and \eqref{eq:abstractinjectivity} is injective exactly when
$\Omega$ is nondegenerate.  No skew form is nondegenerate in odd dimension.  This proves
both necessity and sufficiency.
\end{proof}

\begin{corollary}[Genericity of full rank]\label{cor:genericfullrank}
Assume $b|_{K\times K}=0$ and that the criterion in
Theorem~\ref{thm:abstractfullrank} holds.  Then the full-rank pairings form a nonempty
relative Zariski-open subset of the affine space $\cC(D,b)$; in particular, they are
Euclidean dense.
\end{corollary}

\begin{proof}
In affine coordinates on $\cC(D,b)$, failure of full rank is the common zero set of the
maximal minors of the matrix of $M$.  Its complement is therefore relative Zariski open.
It is nonempty by Theorem~\ref{thm:abstractfullrank}; since an affine space is irreducible,
every nonempty relative Zariski-open subset is Zariski dense, hence Euclidean dense.
\end{proof}

\begin{corollary}[Boundary-free parity obstruction]\label{cor:parityobstruction}
Assume $b=0$.  Then $a_M$ descends to a skew form $\overline a_M$ on $E/K$.  The mixed
pairing has rank $m$ if and only if $\overline a_M$ is nondegenerate.  In particular, a
full-rank compatible pairing exists if and only if $m$ is even.
\end{corollary}

\begin{proof}
For $b=0$, equation \eqref{eq:leftradicalK} puts $K$ in both radicals and $a_M$ descends to
$E/K$.  Since $D^*:F^*\to E^*$ is injective, $M$ and $a_M$ have the same rank.  Thus maximal
rank is equivalent to nondegeneracy of the descended skew form.  Such a form exists exactly
in even dimension.
\end{proof}

The next result isolates the mechanism that will later produce the polynomial endpoint
condition.  For a bilinear form $a$, write
$\Rad_La=\{e\in E:a(e,E)=0\}$.

\begin{theorem}[Radical lines for a rank-two boundary]
\label{thm:abstractradical}
Assume $b|_{K\times K}=0$, $\dim K=1$, $\rank b=2$, and $b(K,E)\neq0$.  Fix
$e\in K\setminus\{0\}$ and put $Z=\Rad b$.  Then the left radical of every full-rank
compatible pairing is a line.  Moreover, the set of all such lines is exactly
\begin{align}
 m\ \text{even}:\quad&
 \bigl\{L\in\Pp(E):L\subset e^{\perp_b},\ L\neq K,\ L\not\subset Z\bigr\},
 \label{eq:abstractradicale}\\[2mm]
 m\ \text{odd}:\quad&
 \bigl\{L=\spanof\{\kappa\}:b(\kappa,\kappa)=0,\ b(e,\kappa)\neq0\bigr\}.
 \label{eq:abstractradicalo}
\end{align}
Every line displayed in \eqref{eq:abstractradicale}--\eqref{eq:abstractradicalo} is realized
constructively by a full-rank compatible pairing.
\end{theorem}

\begin{proof}
Rescale a vector $h$ so that $b(e,h)=2$, and replace $h$ by a multiple of $e$ plus $h$ so
that $b(h,h)=0$.  Then
$E=\spanof\{e,h\}\oplus Z$, $\dim Z=m-1$, and every compatible form has the unique matrix
\begin{equation}\label{eq:abstractranktwonormal}
 A_{r,C}=
 \begin{pmatrix}
 0&2&0\\
 0&0&r^{\mathsf T}\\
 0&-r&C
 \end{pmatrix},
 \qquad C^{\mathsf T}=-C.
\end{equation}
It has rank $m$ exactly when
$\Ker C\cap r^\perp=\{0\}$.  If
$\kappa=\alpha e+\beta h+z$, its left-kernel equations are
\begin{equation}\label{eq:abstractleftkernel}
 z^{\mathsf T}r=2\alpha,
 \qquad Cz=\beta r.
\end{equation}

If $m$ is even, then $\dim Z$ is odd.  Full rank forces
$\Ker C=\spanof\{z\}$ and $z^{\mathsf T}r\neq0$; equation
\eqref{eq:abstractleftkernel} gives $\beta=0$, $z\neq0$, and $\alpha\neq0$.  This is exactly
\eqref{eq:abstractradicale}.  Conversely, for
$\kappa=\alpha e+z$ satisfying those conditions, choose a skew $C$ with kernel
$\spanof\{z\}$ and choose $r$ with $z^{\mathsf T}r=2\alpha$.

If $m$ is odd, then $\dim Z$ is even and full rank forces $C$ to be invertible.  Equations
\eqref{eq:abstractleftkernel} give $\beta\neq0$ and
$\alpha=\tfrac12z^{\mathsf T}r=\tfrac1{2\beta}z^{\mathsf T}Cz=0$.  Since
$b(\kappa,\kappa)=4\alpha\beta$ and $b(e,\kappa)=2\beta$, this is precisely
\eqref{eq:abstractradicalo}.  Conversely, write any such vector as
$\kappa=\beta h+z$, choose an invertible skew $C$ on $Z$, and set
$r=\beta^{-1}Cz$.  In both parities the constructed matrix has rank $m$ and the prescribed
left radical.
\end{proof}

\begin{remark}[Examples and limits of the abstraction]\label{rem:abstractexamples}
The polynomial construction below has $K=\R$ and an endpoint boundary form.  A finite spline
space with $F=DE$ gives the same one-dimensional kernel and boundary trace.  For periodic
trigonometric polynomials, $F=DE$ is the zero-mean space, $b=0$, and $\dim F$ is even, so
Corollary~\ref{cor:parityobstruction} does not forbid full rank.  For graph cochains, the
incidence map is generally not onto the full edge space when cycles are present; the theorem
applies directly to $F=\im D$, while a treatment retaining the cycle space requires the
longer complex and its cohomology.  These examples show both the reach and the precise scope
of the length-one theory.
\end{remark}

\subsection{Scalar localization requires additional moment data}

The quadruple $(E,F,D,b)$ contains no points, evaluation maps, or multiplication.  Locality
therefore cannot be defined from the Stokes structure alone.  Suppose additionally that
$E$ and $F$ are function spaces on a set $X$, and put
\[
 G=\spanof\{ef:e\in E, f\in F\}.
\]
For a linear functional $L\in G^*$, define the product pairing
\begin{equation}\label{eq:abstractmomentpairing}
 M_L(e,f)=L(ef).
\end{equation}
Let
\begin{equation}\label{eq:deltadefinition}
 \delta:S^2E\longrightarrow G,
 \qquad \delta(e,e')=eDe'+e'De.
\end{equation}

\begin{proposition}[Moment-factorization criterion]\label{prop:momentfactorization}
Assume that a reference functional $L_0\in G^*$ satisfies
$L_0(\delta(e,e'))=b(e,e')$.  Then the Stokes-compatible product pairings are precisely
\begin{equation}\label{eq:momentaffine}
 L_0+(\im\delta)^\perp\subset G^*.
\end{equation}
In particular, the product pairing is unique if and only if $\im\delta=G$.
An atomic scalar localization with $r$ nodes and positive weights exists precisely when the
corresponding $L$ has a representation
\begin{equation}\label{eq:atomicmoment}
 L=\sum_{i=1}^r w_i\,\operatorname{ev}_{x_i}|_G,
 \qquad x_i\in X,\quad w_i>0,
\end{equation}
which is a truncated moment problem on $G$.
\end{proposition}

\begin{proof}
Substitution of \eqref{eq:abstractmomentpairing} into \eqref{eq:abstractstokes} gives
$L(\delta(e,e'))=b(e,e')$.  Subtracting the same identity for $L_0$ proves
\eqref{eq:momentaffine}; the remaining statements are immediate.
\end{proof}

\begin{proposition}[Positive localization by Tchakaloff]
\label{prop:tchakaloff}
Let $X$ be compact, let $G\subset C(X)$ be finite dimensional and contain the constants,
and let $L\in G^*$ be representable by a finite positive Borel measure on $X$.  Then there
exist $r\leq\dim G$, nodes $x_i\in X$, and positive weights $w_i$ such that
\begin{equation}\label{eq:tchakaloffrepresentation}
 L(g)=\sum_{i=1}^r w_i g(x_i),\qquad g\in G.
\end{equation}
Consequently, whenever the compatible moment slice is nonempty, its positive localizable
part is
\begin{equation}\label{eq:positivemomentslice}
 \bigl(L_0+(\im\delta)^\perp\bigr)\cap\mathcal M_+(G),
\end{equation}
where $\mathcal M_+(G)$ denotes the cone of functionals represented by positive measures;
every element of this intersection has a positive atomic representative with at most
$\dim G$ nodes.
\end{proposition}

\begin{proof}
Choose a basis $g_1,\ldots,g_s$ of $G$ with $g_1=1$ and consider the continuous moment map
$\Gamma(x)=(g_1(x),\ldots,g_s(x))$.  The moment vector of $L$ lies in the conic hull of
$\Gamma(X)$.  The conic Carath\'eodory theorem represents it as a positive combination of at
most $s=\dim G$ such vectors, giving \eqref{eq:tchakaloffrepresentation}.  Formula
\eqref{eq:positivemomentslice} then follows from
Proposition~\ref{prop:momentfactorization}.  This is Tchakaloff's theorem in the present
finite-dimensional form \cite{tchakaloff1957,bayerteichmann2006}.
\end{proof}

Complete Chebyshev systems provide classical hypotheses under which generalized Gaussian
representations exist and are unique; weak Chebyshev systems lead to related moment and
spline quadrature theories \cite{karlinstudden1966,gustafson1970,micchellipinkus1977,cheng1999}.
These results concern the additional moment cone in \eqref{eq:atomicmoment}, not the bare
Stokes quadruple.  Consequently, a classification of all function spaces admitting a
minimal positive localization is a separate moment-theoretic problem rather than a corollary
of Theorem~\ref{thm:abstractaffine}.

Thus three levels must be kept separate: affine compatibility of a mixed pairing, positive
representability of a moment functional, and existence or uniqueness of a finite atomic
representation.  Tchakaloff shows that every positively representable functional admits a
positive atomic representation with at most $\dim G$ atoms.  It neither proves that the
compatible affine slice meets the positive moment cone nor guarantees a minimal node count;
Chebyshev-system theory supplies sharper counts and uniqueness only under additional
structure.

\section{Polynomial mixed Stokes pairings}

\subsection{Spaces, derivative, and boundary form}

Fix $N\geq1$ and set
\[
 E_N=\Pp_N[-1,1],\qquad F_N=\Pp_{N-1}[-1,1].
\]
Throughout the polynomial sections, $P_n$ denotes the Legendre polynomial normalized by
$P_n(1)=1$.
The derivative
\[
 d:E_N\longrightarrow F_N,\qquad d q=q',
\]
is surjective and has kernel $\R$.  Define the symmetric boundary form
\begin{equation}\label{eq:boundaryform}
 b(p,q)=\bracket{pq}=p(1)q(1)-p(-1)q(-1).
\end{equation}

\begin{definition}[Stokes-compatible mixed pairing]
A bilinear map $\beta:E_N\times F_N\to\R$ is \emph{Stokes-compatible} if
\begin{equation}\label{eq:stokesbeta}
 \beta(p,q')+\beta(q,p')=b(p,q),
 \qquad p,q\in E_N.
\end{equation}
It is \emph{full rank} if the map
\[
 \beta^\flat:E_N\to F_N^*,\qquad
 p\mapsto\beta(p,\cdot),
\]
has rank $N$.
\end{definition}

In bases of $E_N$ and $F_N$, let $M\in\R^{(N+1)\times N}$ represent $\beta$ and
$D\in\R^{N\times(N+1)}$ represent $d$.  If $R$ evaluates at the endpoints and
$\Sigma=\diag(-1,1)$, then \eqref{eq:stokesbeta} is exactly
\begin{equation}\label{eq:matrixstokes}
 MD+D^{\mathsf T}M^{\mathsf T}=R^{\mathsf T}\Sigma R.
\end{equation}
This is the rectangular discrete Stokes identity of the 2012 construction
\cite{moulla2012}.

The reference pairing is
\begin{equation}\label{eq:l2pairing}
 \beta_0(p,r)=\int_{-1}^{1}p(x)r(x)\dd x.
\end{equation}
Integration by parts shows that $\beta_0$ is Stokes-compatible.

\subsection{From a rectangular pairing to a bilinear form on efforts}

Associate with $\beta$ the form
\begin{equation}\label{eq:abeta}
 a_\beta(p,q)=\beta(p,q'),\qquad p,q\in E_N.
\end{equation}
It satisfies
\begin{equation}\label{eq:asympart}
 a_\beta(p,q)+a_\beta(q,p)=b(p,q),
 \qquad a_\beta(p,1)=0.
\end{equation}
Conversely, every bilinear form satisfying \eqref{eq:asympart} determines a unique mixed
pairing.  Indeed, for $r\in F_N$, choose any primitive $q\in E_N$ with $q'=r$ and set
$\beta(p,r)=a_\beta(p,q)$.  This is independent of the additive constant in $q$.

For the reference pairing,
\[
 a_0(p,q)=\int_{-1}^{1}p(x)q'(x)\dd x.
\]
The symmetric part of $a_0$ is the boundary form, not the $L^2$ product.  The role of
$L^2$ is that of a distinguished reference point in the affine class.

\section{Affine classification and pairing gauge}

\begin{theorem}[Affine classification]\label{thm:affine}
The set $\cA_N$ of Stokes-compatible mixed pairings is an affine space with origin
$\beta_0$ and model vector space
\[
 \Lambda^2(E_N/\R)^*.
\]
Equivalently, every compatible form has a unique decomposition
\begin{equation}\label{eq:gaugedecomposition}
 a_\beta=a_0+\omega,
 \qquad
 \omega\in\Lambda^2(E_N/\R)^*.
\end{equation}
Consequently,
\begin{equation}\label{eq:gaugedimension}
 \dim\cA_N=\frac{N(N-1)}2.
\end{equation}
\end{theorem}

\begin{proof}
Let $\beta_1$ and $\beta_0$ be compatible and put
$\omega=a_{\beta_1}-a_{\beta_0}$.  Subtracting their Stokes identities gives
\[
 \omega(p,q)+\omega(q,p)=0.
\]
Thus $\omega$ is skew-symmetric.  Since both forms vanish in the second slot on constants,
$\omega(p,1)=0$; skew-symmetry also gives $\omega(1,p)=0$.  Hence $\omega$ descends to a
two-form on $E_N/\R$.

Conversely, if $\omega\in\Lambda^2(E_N/\R)^*$, then $a_0+\omega$ vanishes on constants in
the second slot and has symmetric part $b$.  It therefore defines a unique compatible mixed
pairing.  Since $\dim(E_N/\R)=N$, the dimension is $N(N-1)/2$.
\end{proof}

\begin{definition}[Pairing gauge]
The two-form $\omega$ in \eqref{eq:gaugedecomposition} is called the \emph{pairing gauge}.
The reference $\omega=0$ is the \emph{$L^2$ gauge}.
\end{definition}

The terminology records a power-invisible antisymmetric freedom.  It is not a basis change:
different $\omega$ define different bilinear pairings.  Since $E_N/\R$ is a vector space and
$\omega$ is constant, it is also a closed two-form, which anticipates the Dirac-geometric
interpretation in Section~\ref{sec:continuum}.

In matrix form, if $M_0$ represents $\beta_0$ and $\Delta M=M-M_0$, then
\begin{equation}\label{eq:omegamatrix}
 \Omega=\Delta M D
\end{equation}
represents $\omega$.  Equation \eqref{eq:matrixstokes} gives
$\Omega^{\mathsf T}=-\Omega$, and $D1=0$ gives $\Omega1=0$.

\section{Canonical normal form and the full-rank locus}

Assume $N\geq2$.  Introduce
\[
 e=1,\qquad h=x,\qquad
 Z_N=(1-x^2)\Pp_{N-2}.
\]
Then
\begin{equation}\label{eq:decompositionE}
 E_N=\spanof\{e,h\}\oplus Z_N,
 \qquad d_Z:=\dim Z_N=N-1.
\end{equation}
Every element is written uniquely as
$p=\alpha e+\beta h+z$, with $z\in Z_N$.  In this decomposition, the boundary form has
matrix
\begin{equation}\label{eq:Bnormal}
 [b]=
 \begin{pmatrix}
 0&2&0\\
 2&0&0\\
 0&0&0
 \end{pmatrix}.
\end{equation}

\begin{proposition}[Canonical normal form]\label{prop:normalform}
Every Stokes-compatible form $a_\beta$ has a unique matrix
\begin{equation}\label{eq:Anormal}
 A_{r,C}=
 \begin{pmatrix}
 0&2&0\\
 0&0&r^{\mathsf T}\\
 0&-r&C
 \end{pmatrix},
 \qquad r\in\R^{d_Z},\quad C\in\R^{d_Z\times d_Z},\quad C^{\mathsf T}=-C.
\end{equation}
Conversely, every matrix \eqref{eq:Anormal} defines a Stokes-compatible pairing.
Moreover,
\begin{equation}\label{eq:rankcondition}
 \rank A_{r,C}=N
 \quad\Longleftrightarrow\quad
 \Ker C\cap r^\perp=\{0\}.
\end{equation}
\end{proposition}

\begin{proof}
The first column of $A$ is zero because $a_\beta(\cdot,1)=0$.  The equation
$A+A^{\mathsf T}=[b]$ then fixes the first row, the $2\times2$ block, and the off-diagonal
signs, leaving precisely a vector $r$ and a skew matrix $C$.

For $u=u_e e+u_hh+u_z$, the equation $A_{r,C}u=0$ gives
\[
 u_h=0,\qquad Cu_z=0,\qquad r^{\mathsf T}u_z=0.
\]
Thus the right kernel is
$\spanof\{e\}\oplus(\Ker C\cap r^\perp)$.  Since the maximal possible rank is $N$, the
rank criterion follows.
\end{proof}

Because $D$ is surjective, $\rank(MD)=\rank M$.  Hence full rank of $\beta$ is equivalent to
$\rank A_{r,C}=N$.  For a full-rank pairing the left radical is one-dimensional.  If
\begin{equation}\label{eq:kappadecomp}
 \kappa=\alpha e+\beta h+z,\qquad z\in Z_N,
\end{equation}
then $\kappa^{\mathsf T}A_{r,C}=0$ is equivalent to
\begin{equation}\label{eq:leftkerneleq}
 z^{\mathsf T}r=2\alpha,
 \qquad
 Cz=\beta r.
\end{equation}

\section{The kernel map and its exact image}

Let $\cA_N^{\mathrm{fr}}\subset\cA_N$ denote the open subset of full-rank pairings.  Define
the projective kernel map
\begin{equation}\label{eq:kernelmap}
 \cK_N:\cA_N^{\mathrm{fr}}\longrightarrow\Pp(E_N),
 \qquad
 \beta\longmapsto[\Ker\beta^\flat]=[\kappa_\beta].
\end{equation}

\begin{theorem}[Exact image of the kernel map]\label{thm:kernelimage}
For $N\geq1$, the image of $\cK_N$ consists exactly of the projective classes of nonconstant
polynomials $\kappa\in E_N$ satisfying
\begin{equation}\label{eq:endpointparity}
 \boxed{\kappa(1)=(-1)^N\kappa(-1)\neq0.}
\end{equation}
In particular, the degree-$N$ stratum of the image is precisely the set of degree-$N$
polynomials satisfying \eqref{eq:endpointparity}.
If
\begin{equation}\label{eq:projectivehyperplane}
 \mathcal H_N=
 \Pp\!\left(\Ker\bigl[\kappa\mapsto
 \kappa(1)-(-1)^N\kappa(-1)\bigr]\right),
\end{equation}
then $\im\cK_N$ is the relative semialgebraic open subset of $\mathcal H_N$ obtained by
removing the endpoint-vanishing locus and, when $N$ is even, the constant line.
\end{theorem}

\begin{proof}
Here $K=\R e$, $m=N$, $\rank b=2$, and
$b(e,p)=p(1)-p(-1)$.  If $N$ is even,
Theorem~\ref{thm:abstractradical} says that the radical line lies in $e^{\perp_b}$, is not
the constant line, and is not contained in $\Rad b=Z_N$.  These three conditions are
respectively
\[
 \kappa(1)=\kappa(-1),\qquad \kappa\ \text{nonconstant},\qquad
 \kappa(1)=\kappa(-1)\neq0.
\]
If $N$ is odd, the abstract conditions become
\[
 \kappa(1)^2-\kappa(-1)^2=0,\qquad
 \kappa(1)-\kappa(-1)\neq0,
\]
which are equivalent to
$\kappa(1)=-\kappa(-1)\neq0$.  The constructive converse in
Theorem~\ref{thm:abstractradical} proves sufficiency in both parities.  Finally, the excluded
conditions are algebraic closed subsets of the projective hyperplane
\eqref{eq:projectivehyperplane}; their complement is therefore relative semialgebraic open.
\end{proof}

\begin{remark}[Degree drop]
Full rank does not force $\deg\kappa_\beta=N$.  Lower-degree kernel polynomials occur on
proper algebraic strata of the image.  The node theorem below therefore concerns the open
degree-$N$ stratum.  This distinction is essential: only a degree-$N$ kernel can supply $N$
finite roots counted with multiplicity.
\end{remark}

The endpoint value in Theorem~\ref{thm:kernelimage} gives every projective class in the
image a unique representative satisfying $\kappa(1)=1$.  This normalization reveals the
topology hidden by the projective notation.

\begin{theorem}[Topology of the kernel image]\label{thm:basetopology}
Let $\cB_N=\im\cK_N$, represented by the normalized polynomials $\kappa(1)=1$.  Then
\begin{equation}\label{eq:basetopology}
 \boxed{
 \cB_N\simeq
 \begin{cases}
  \R^{N-1},&N\ \text{odd},\\
  \R^{N-1}\setminus\{0\},&N\ \text{even}.
 \end{cases}}
\end{equation}
More precisely, for $N\geq2$, with $Z_N=(1-x^2)\Pp_{N-2}$,
\begin{equation}\label{eq:basenormalized}
 \kappa=x+z\quad(N\ \text{odd}),
 \qquad
 \kappa=1+z,\ z\neq0\quad(N\ \text{even}).
\end{equation}
Thus $\cB_N$ is contractible for odd $N$; for $N=2$ it has two components; and for even
$N\geq4$ it is connected and deformation retracts onto $S^{N-2}$.

For $N\geq2$, the exact-degree stratum
$\cB_N^{(N)}=\{[\kappa]\in\cB_N:\deg\kappa=N\}$ has exactly two contractible components,
distinguished projectively by the sign of
\begin{equation}\label{eq:leadinvariant}
 c_N(\kappa)\kappa(1),
\end{equation}
where $c_N(\kappa)$ is the coefficient of $P_N$ in the Legendre expansion of $\kappa$.
The chamber of polynomials with $N$ simple roots in $(-1,1)$ lies in the positive component.
For $N=1$, the base and its exact-degree stratum are the single class $[x]$.
\end{theorem}

\begin{proof}
For $N=1$ the base is $[x]$.  Assume $N\geq2$.  In the decomposition
$E_N=\spanof\{1,x\}\oplus Z_N$, the normalized endpoint condition
forces respectively the two representations in \eqref{eq:basenormalized}.  Since
$\dim Z_N=N-1$, this proves \eqref{eq:basetopology} and its stated homotopy consequences.

For $N\geq2$, exact degree is the complement of the affine hyperplane
$c_N(\kappa)=0$ in these coordinates.  Its two open half-spaces are contractible.  Under a
projective rescaling both factors in \eqref{eq:leadinvariant} acquire the same scalar, so
its sign is intrinsic.  If $\kappa$ has simple interior roots, its monic representative has
positive leading Legendre coefficient and positive value at $1$; hence the sign is positive.
This completes the proof.
\end{proof}

\subsection{Real node configurations}

Let $\kappa$ be monic of degree $N$ with roots $x_1,\ldots,x_N$.  Then
\[
 \kappa(1)=\prod_{i=1}^N(1-x_i),
 \qquad
 (-1)^N\kappa(-1)=\prod_{i=1}^N(1+x_i).
\]
Theorem~\ref{thm:kernelimage} immediately gives the following classification.

\begin{corollary}[Admissible interior nodes]\label{cor:nodes}
Let $-1<x_1<\cdots<x_N<1$.  There exists a full-rank Stokes-compatible pairing whose
kernel polynomial has precisely these roots if and only if
\begin{equation}\label{eq:productbalance}
 \boxed{\prod_{i=1}^N(1-x_i)=\prod_{i=1}^N(1+x_i).}
\end{equation}
Equivalently,
\begin{equation}\label{eq:hyperbolicbalance}
 \boxed{\sum_{i=1}^N\operatorname{artanh}(x_i)=0.}
\end{equation}
\end{corollary}

\begin{proof}
The product identity is the endpoint condition.  Since all factors are positive, taking
logarithms and using
$2\operatorname{artanh}x=\log((1+x)/(1-x))$ gives \eqref{eq:hyperbolicbalance}.
Sufficiency follows from Theorem~\ref{thm:kernelimage}.
\end{proof}

Let
\[
 \mathfrak C_N=\{(x_1,\ldots,x_N):-1<x_1<\cdots<x_N<1\}.
\]
The map $\Phi(x)=\sum_i\operatorname{artanh}(x_i)$ has nonzero gradient
\[
 \nabla\Phi(x)=\left(\frac1{1-x_1^2},\ldots,\frac1{1-x_N^2}\right).
\]
Hence the admissible ordered node configurations form the smooth hypersurface
$\Phi^{-1}(0)\subset\mathfrak C_N$, of dimension $N-1$.  Stokes therefore imposes a genuine
codimension-one obstruction: it does not generate arbitrary grids.

\begin{proposition}[Hyperbolic barycentre and configuration geometry]
\label{prop:hyperbolicgeometry}
Equip $I=(-1,1)$ with the one-dimensional Klein--Hilbert metric
\begin{equation}\label{eq:hyperbolicmetric}
 g_I=\frac{\dd x^2}{(1-x^2)^2}.
\end{equation}
Then $y=\operatorname{artanh}x$ is a global geodesic coordinate and an isometry from
$(I,g_I)$ to the Euclidean line.  For $X=(x_1,\ldots,x_N)\in\mathfrak C_N$, condition
\eqref{eq:hyperbolicbalance} is equivalent to saying that the equal-weight Fr\'echet
barycentre of $x_1,\ldots,x_N$ is the origin.  In product rapidity coordinates, the
admissible configuration space is
\begin{equation}\label{eq:rapiditychamber}
 \left\{y_1<\cdots<y_N:\ \sum_{i=1}^Ny_i=0\right\},
\end{equation}
an open chamber in a totally geodesic hyperplane of $\R^N$.
\end{proposition}

\begin{proof}
Equation \eqref{eq:hyperbolicmetric} gives $\dd s=\dd x/(1-x^2)=\dd y$.  Hence the
Fr\'echet functional of a candidate barycentre with rapidity $z$ is
$\sum_i(z-y_i)^2$, whose unique minimizer is $z=N^{-1}\sum_i y_i$.  It is the origin
exactly under \eqref{eq:hyperbolicbalance}.  The product metric becomes Euclidean in the
$y_i$, proving the last statement.
\end{proof}

\begin{remark}[What the hyperbolic interpretation does not assert]
The endpoints $\pm1$ are the two ideal points of the interval model, and the Stokes form
uses traces at those endpoints.  This observation does not by itself make the polynomial
Stokes structure equivariant under the hyperbolic isometry group: polynomial spaces are not
preserved by general M\"obius reparametrizations.  Nor does
Proposition~\ref{prop:hyperbolicgeometry} produce a multivariate Gaussian cubature on the
Poincar\'e disk.  Such an extension would require a separate differential complex, boundary
trace, moment space, and zero-set theory.
\end{remark}

\subsection{Fibres of the kernel map}

\begin{theorem}[Fibre dimension]\label{thm:fibres}
For every projective kernel class in the image of $\cK_N$, the corresponding full-rank fibre
is a nonempty smooth open subset of an affine space of dimension
\begin{equation}\label{eq:fibredim}
 \boxed{\frac{(N-1)(N-2)}2.}
\end{equation}
\end{theorem}

\begin{proof}
Use the decomposition $\kappa=\alpha e+\beta h+z$.

For even $N$, $\beta=0$, $z\neq0$.  The condition $Cz=0$ leaves a skew form on the
$(d_Z-1)$-dimensional quotient $Z_N/\spanof\{z\}$, of dimension
$(d_Z-1)(d_Z-2)/2$.  The condition $z^{\mathsf T}r=2\alpha$ leaves $d_Z-1$ free components of
$r$.  The total is $d_Z(d_Z-1)/2=(N-1)(N-2)/2$.  Full rank is the open condition that the
induced skew form on the even-dimensional quotient be nondegenerate.

For odd $N$, $C$ may be any invertible skew matrix and $r=\beta^{-1}Cz$ is then fixed.
Invertible skew matrices form an open subset of dimension $d_Z(d_Z-1)/2$, giving the same
formula.
\end{proof}

The dimension identity
\[
 \frac{N(N-1)}2
 = (N-1)+\frac{(N-1)(N-2)}2
\]
separates the $N-1$ observable kernel directions from the gauge directions invisible to the
kernel map.

\begin{theorem}[Topology of the full-rank fibres]\label{thm:fibretopology}
For $N=1$ and $N=2$, every nonempty full-rank fibre is a point.  For every $N\geq3$, each
full-rank fibre has exactly two connected components.  After orienting the relevant even
dimensional quotient, the components are distinguished by the sign of a Pfaffian.
\end{theorem}

\begin{proof}
For odd $N\geq3$, the proof of Theorem~\ref{thm:fibres} identifies the fibre with the space
of invertible skew forms $C$ on the even-dimensional space $Z_N$; $r$ is then fixed by $C$.
The sign of $\operatorname{Pf}(C)$ is constant on connected components, and the two sign
sets are connected because each is a homogeneous space of one component of
$\operatorname{GL}(Z_N)$ by the symplectic group.

For even $N$, the prescribed kernel fixes $z\neq0$.  The skew form $C$ descends to a
nondegenerate skew form on the even-dimensional quotient
$Z_N/\spanof\{z\}$, while $r$ ranges over the affine hyperplane
$z^{\mathsf T}r=2\alpha$.  The latter is contractible, so the two components are again the
two Pfaffian signs of the quotient form.  When $N=2$ that quotient is zero dimensional and
the affine condition fixes $r$, leaving a point.  For $N=1$ the gauge space itself is zero
dimensional.
\end{proof}

Pointwise synthesis in Theorem~\ref{thm:kernelimage} does not automatically give a continuous
choice of pairing over the whole base.  The topology of Theorem~\ref{thm:basetopology}
produces an exact obstruction.

\begin{theorem}[Global continuous and smooth full-rank sections]
\label{thm:globalsmoothsection}
The projective kernel map
\begin{equation}\label{eq:sectionmap}
 \cK_N:\cA_N^{\mathrm{fr}}\longrightarrow\cB_N
\end{equation}
admits a global continuous section if and only if it admits a global smooth section, and
these equivalent conditions hold if and only if
\begin{equation}\label{eq:sectionorders}
 \boxed{N\ \text{is odd}\quad\text{or}\quad N\in\{2,4,8\}.}
\end{equation}
For even $N$, the obstruction is equivalent to the existence of an almost complex structure
on $S^{N-2}$.
\end{theorem}

\begin{proof}
For $N=1$ the assertion is immediate.  Suppose first that $N\geq3$ is odd.  Then
$d_Z=N-1$ is even.  Fix a nondegenerate skew matrix $C_0$ on $Z_N$.  For the normalized
kernel $\kappa=x+z$, set
\begin{equation}\label{eq:oddglobalsection}
 C_z=C_0,\qquad r_z=C_0z.
\end{equation}
The equations \eqref{eq:leftkerneleq} hold, $C_0$ is invertible, and
Proposition~\ref{prop:normalform} gives a full-rank pairing depending smoothly on $z$.

Let now $N$ be even.  Write the normalized kernel as $\kappa=1+z$ with
$z\in Z_N\setminus\{0\}$.  Equations \eqref{eq:leftkerneleq} and the full-rank condition are
equivalent to
\begin{equation}\label{eq:evensectiondata}
 C_z^{\mathsf T}=-C_z,\qquad \Ker C_z=\spanof\{z\},
 \qquad z^{\mathsf T}r_z=2.
\end{equation}
Fix an auxiliary Euclidean inner product on $Z_N$, use its musical isomorphism
$z\mapsto z^\flat$, and set
\begin{equation}\label{eq:globalrz}
 r_z=\frac{2z^\flat}{\|z\|^2}.
\end{equation}
Thus all topology is carried by $C_z$.  Restricting a family $C_z$ to the unit sphere in
$Z_N$ gives the tangent two-form
\begin{equation}\label{eq:tangenttwoform}
 \omega_z(v,w)=\langle v,C_zw\rangle,
 \qquad v,w\in T_zS(Z_N)=z^\perp.
\end{equation}
It is nondegenerate precisely because $\Ker C_z=\spanof\{z\}$.  Hence a continuous section
produces a continuous almost symplectic structure on
\begin{equation}\label{eq:tangentsphere}
 T_zS(Z_N)=z^\perp.
\end{equation}
Conversely, an almost symplectic structure, together with the auxiliary metric, produces the
required skew endomorphism $C_z$, extended by zero on $\spanof\{z\}$ and radially on
$Z_N\setminus\{0\}$.  To justify smoothing, use the auxiliary metric to realize
$\Lambda^2T^*S^{N-2}$ as a smooth vector subbundle of the trivial bundle
$S^{N-2}\times\Lambda^2Z_N^*$.  Approximate a continuous section as an ambient map by the
relative smooth approximation theorem \cite[Theorem~2.2.5]{hirsch1976}, then apply the
smooth fibrewise orthogonal projection back to $\Lambda^2T^*S^{N-2}$.  On the compact
sphere, a continuous nondegenerate form has a positive uniform lower bound on its least
tangent singular value, so every sufficiently close smooth approximation remains
nondegenerate.  Thus continuous and smooth existence are equivalent.

An almost symplectic structure is equivalent to an almost complex structure.  For $N=2$
the tangent bundle of $S^0$ is zero dimensional.  Among positive-dimensional spheres, only
$S^2$ and $S^6$ admit almost complex structures
\cite{borelserre1953,konstantisparton2017}.  Hence the even orders are exactly $N=2,4,8$.
\end{proof}

The obstruction concerns the whole projective kernel image, including its lower-degree
strata.  It disappears on the exact-degree set $\cB_N^{(N)}$: its connected components are
contractible, and Theorem~\ref{thm:globalcanonicalrank} below provides a canonical full-rank
section.  Thus the obstruction does not restrict the admissible simple node configurations.

\begin{proposition}[Explicit exceptional sections]\label{prop:exceptionalsections}
For $N=2,4,8$, the even-order sections in Theorem~\ref{thm:globalsmoothsection} may be chosen
explicitly.  With the auxiliary metric fixed as above, write
$\widehat z=z/\lVert z\rVert$ for $z\ne0$:
\begin{enumerate}
\item for $N=2$, take $C_z=0$;
\item for $N=4$, identify the oriented space $Z_4$ with $\R^3$ and set
      $C_zv=\widehat z\times v$;
\item for $N=8$, identify $Z_8$ with the imaginary octonions and use the seven-dimensional
      vector cross product $C_zv=\widehat z\times v$.
\end{enumerate}
In every case take $r_z=2z^\flat/\|z\|^2$.
\end{proposition}

\begin{proof}
In dimensions three and seven, the indicated cross products satisfy
$\langle \widehat z\times v,w\rangle=-\langle v,\widehat z\times w\rangle$ and
$\Ker(v\mapsto\widehat z\times v)=\spanof\{z\}$ for $z\neq0$; see
\cite{browngray1967,baez2002}.  Formula \eqref{eq:globalrz} supplies the remaining normal-form
condition.  Moreover, $C_z^2=-(I-\widehat z\otimes\widehat z^\flat)$, so the nonzero tangent
singular values equal one.  The case $N=2$ has $\dim Z_2=1$ and zero-dimensional tangent
spaces.
\end{proof}

We call a \emph{synthesis chart} an open subset of $\cB_N$ endowed with a smooth full-rank
local section of $\cK_N$.

\begin{corollary}[Minimal synthesis atlas]\label{cor:minimalatlas}
The minimum number of smooth full-rank synthesis charts covering $\cB_N$ is one when $N$ is
odd or $N\in\{2,4,8\}$, and two for every other even $N$.
\end{corollary}

\begin{proof}
One chart exists exactly in the orders of Theorem~\ref{thm:globalsmoothsection}.  In the
remaining even orders, one chart is impossible.  Two suffice because the sphere
$S^{N-2}$ is covered by two punctured spheres; each is contractible, so its tangent bundle
admits a smooth symplectic frame.  Extend the corresponding nondegenerate tangent form
radially and use \eqref{eq:globalrz}.
\end{proof}

\begin{remark}[Reproducible auxiliary choices]\label{rem:reproduciblechoices}
The atlas becomes deterministic once auxiliary conventions are fixed.  One convenient choice
is the $L^2$ metric on $Z_N$, followed by Gram--Schmidt orthonormalization of the ordered
basis
\[
 (1-x^2),\ x(1-x^2),\ldots,x^{N-2}(1-x^2),
\]
with positive normalization factors.  Write the resulting oriented basis as
$(e_1,\ldots,e_{N-1})$, take $e=e_1$, set $J_0e_1=0$, and define
\[
 J_0e_{2j}=e_{2j+1},\qquad
 J_0e_{2j+1}=-e_{2j},
 \qquad 1\leq j\leq\frac{N-2}{2}.
\]
These conventions are computational rather than intrinsic; changing them changes the local
section but not the atlas size, the kernel line, or the Stokes identity.
\end{remark}

\begin{proposition}[Explicit two-chart synthesis]\label{prop:explicittwocharts}
Let $N$ be even.  Fix a unit vector $e\in Z_N$ and a skew endomorphism $J_0$ such that
$\Ker J_0=\spanof\{e\}$ and $J_0^2=-I$ on $e^\perp$.  Such a $J_0$ exists because
$\dim e^\perp=N-2$ is even.  For $z\neq0$ write
$\widehat z=z/\|z\|$.  Given $\varepsilon\in\{+1,-1\}$, put
\begin{equation}\label{eq:rodriguesdata}
 e_\varepsilon=\varepsilon e,\qquad
 U_\varepsilon=\{\widehat z\in S(Z_N):\widehat z\neq-e_\varepsilon\},
\end{equation}
and, on $U_\varepsilon$,
\begin{equation}\label{eq:rodriguesQ}
 A_\varepsilon(\widehat z)
 =\widehat z\otimes e_\varepsilon^\flat
   -e_\varepsilon\otimes\widehat z^\flat,
 \qquad
 Q_\varepsilon(\widehat z)
 =I+A_\varepsilon(\widehat z)
   +\frac{A_\varepsilon(\widehat z)^2}
          {1+\langle e_\varepsilon,\widehat z\rangle}.
\end{equation}
Then
\begin{equation}\label{eq:explicittwocharts}
 C_z^{(\varepsilon)}
 =Q_\varepsilon(\widehat z)J_0Q_\varepsilon(\widehat z)^*,
 \qquad
 r_z=\frac{2z^\flat}{\|z\|^2}
\end{equation}
defines a smooth full-rank synthesis on the cone over $U_\varepsilon$.  The two cones cover
$Z_N\setminus\{0\}$ and therefore give an explicit two-chart atlas at every even order.
\end{proposition}

\begin{proof}
The Rodrigues identity in \eqref{eq:rodriguesQ} gives
$Q_\varepsilon^*Q_\varepsilon=I$ and
$Q_\varepsilon e_\varepsilon=\widehat z$.  Since
$\spanof\{e_\varepsilon\}=\spanof\{e\}$, conjugation yields
$(C_z^{(\varepsilon)})^*=-C_z^{(\varepsilon)}$ and
$\Ker C_z^{(\varepsilon)}=\spanof\{\widehat z\}=\spanof\{z\}$.
Together with $r_z(z)=2$, these are exactly the full-rank conditions
\eqref{eq:evensectiondata}.  Finally $U_+\cup U_-=S(Z_N)$.
\end{proof}

\begin{corollary}[Metric normalization and stable chart selection]
\label{cor:chartnormalization}
For either chart in Proposition~\ref{prop:explicittwocharts},
\begin{equation}\label{eq:Cnormalization}
 \boxed{
 (C_z^{(\varepsilon)})^2
 =-\bigl(I-\widehat z\otimes\widehat z^\flat\bigr).}
\end{equation}
Hence $C_z^{(\varepsilon)}$ vanishes on $\spanof\{z\}$, restricts to an orthogonal complex
structure on $z^\perp$, and has all nonzero singular values equal to one.  Moreover,
\begin{equation}\label{eq:rminimum}
 \boxed{
 \|r_z\|=\frac{2}{\|z\|}
 =\min\{\|r\|:r\in Z_N^*,\ r(z)=2\},}
\end{equation}
and $r_z$ is the unique minimizer.  At evaluation time, choose
$\varepsilon\in\{+1,-1\}$ so that
$\varepsilon\langle e,\widehat z\rangle\geq0$; equivalently, take the sign of the inner
product when it is nonzero and set $\varepsilon=+1$ when it vanishes.  Then the Rodrigues
denominator satisfies
\begin{equation}\label{eq:stablechartdenominator}
 1+\langle e_\varepsilon,\widehat z\rangle
 =1+|\langle e,\widehat z\rangle|\in[1,2].
\end{equation}
\end{corollary}

\begin{proof}
Since $e$ is unit and $J_0$ is an orthogonal complex structure on $e^\perp$,
$J_0^2=-(I-e\otimes e^\flat)$.  Orthogonal conjugation by $Q_\varepsilon$ and
$Q_\varepsilon e=\varepsilon\widehat z$ give \eqref{eq:Cnormalization}.  For every covector
$r$ with $r(z)=2$, Cauchy--Schwarz gives $2\leq\|r\|\|z\|$, with equality exactly when
$r$ is a positive multiple of $z^\flat$.  The constraint then fixes that multiple and gives
\eqref{eq:rminimum}.  The final identity follows directly from the stated choice of sign.
\end{proof}

\begin{corollary}[Unavoidable radial blow-up]\label{cor:radialblowup}
Let $N$ be even and normalize an admissible kernel as $\kappa=1+z$, with
$z\in Z_N\setminus\{0\}$.  In the fixed Euclidean normal-form coordinates, every compatible
pairing having $\kappa$ in its left radical has data $(r_z,C_z)$ satisfying
\begin{equation}\label{eq:unavoidable-r-bound}
 \boxed{\|r_z\|\geq\frac{2}{\|z\|}.}
\end{equation}
Equality holds only for $r_z=2z^\flat/\|z\|^2$.  If $A_z=A_{r_z,C_z}$ is the composed
Stokes matrix in \eqref{eq:Anormal}, then
\begin{equation}\label{eq:unavoidable-A-bound}
 \boxed{\|A_z\|_{\mathrm F}\geq\frac{2\sqrt2}{\|z\|}.}
\end{equation}
If $M_zD=A_z$ in fixed orthonormal coefficient coordinates, then
\begin{equation}\label{eq:unavoidable-M-bound}
 \boxed{\|M_z\|_{\mathrm F}
 \geq\frac{2\sqrt2}{\|D\|_2\,\|z\|}.}
\end{equation}
Consequently, no full-rank synthesis can remain bounded while approaching the deleted
constant class $z=0$, including in the exceptional orders $N=2,4,8$.
The numerical constants depend on the auxiliary metric and fixed coefficient coordinates,
but the unboundedness and the lower-bound order $\|z\|^{-1}$ are invariant, up to
multiplicative constants, under every fixed change of coordinates.
\end{corollary}

\begin{proof}
The general radical equation \eqref{eq:leftkerneleq} gives $r_z(z)=2$, so Cauchy--Schwarz yields
\eqref{eq:unavoidable-r-bound} and its equality case.  In the normal form
\eqref{eq:Anormal}, the vector $r_z$ occurs in two disjoint blocks; hence
\[
 \|A_z\|_{\mathrm F}^2
 =4+2\|r_z\|^2+\|C_z\|_{\mathrm{HS}}^2
 \geq2\|r_z\|^2,
\]
which proves \eqref{eq:unavoidable-A-bound}.  Finally,
$\|M_zD\|_{\mathrm F}\leq\|M_z\|_{\mathrm F}\|D\|_2$ gives
\eqref{eq:unavoidable-M-bound}.  For the explicit normalized sections,
$\|C_z\|_{\mathrm{HS}}^2=N-2$ and $\|r_z\|^2=4/\|z\|^2$, so the divergence rate
$\|z\|^{-1}$ is attained by $A_z$.  Since $D$ is fixed and surjective, the reconstruction
$A\mapsto M$ defined by $MD=A$ is a linear isomorphism from the forms vanishing on $\Ker D$
in the second slot to the mixed pairings, and both it and its inverse are bounded.  Hence the
same rate is attained by $M_z$, up to coordinate-dependent constants.  Only the radial
component is norm-minimal; the unit tangent normalization is a conditioning convention, not a
minimum-norm claim for $C_z$.
\end{proof}

\begin{remark}[Algorithmic meaning]
Every admissible kernel has a full-rank realization at every order.  Proposition
\ref{prop:explicittwocharts} turns this pointwise fact into a closed-form, optimization-free
synthesis with two explicit charts: the formulas give $(C_z,r_z)$, followed by the fixed
linear reconstruction $M_zD=A_{r_z,C_z}$.  Corollary~\ref{cor:chartnormalization} gives a
uniformly normalized tangent block and an antipode-safe chart selector.  It does not give a
uniformly bounded full-rank synthesis near the deleted constant line, where
Corollary~\ref{cor:radialblowup} forces $\|r_z\|\geq2/\|z\|$.
Theorem~\ref{thm:globalsmoothsection} says exactly when those charts can be replaced by one
global smooth rule.  The sign selector is piecewise smooth and may jump across the equator;
in the nonexceptional even orders this is consistent with the obstruction to a global
section.  The obstruction belongs to the choice of gauge over the base, not to the existence
of individual fibres.
\end{remark}

\section{Low-order models}

\subsection{The complete \texorpdfstring{$N=2$}{N=2} family}

Use the monomial bases $(1,x,x^2)$ of $E_2$ and $(1,x)$ of $F_2$.  Every compatible pairing
is represented, after the normalization used below, by
\begin{equation}\label{eq:M2}
 M_t=
 \begin{pmatrix}
 2&0\\
 0&\frac23+\frac t2\\
 \frac23-t&0
 \end{pmatrix}.
\end{equation}
Its kernel polynomial is
\begin{equation}\label{eq:kappa2}
 \kappa_t(x)=x^2+\frac t2-\frac13.
\end{equation}
When
\[
 a^2=\frac13-\frac t2\in(0,1),
\]
the intrinsic nodes are $\{-a,a\}$.  Thus the one-dimensional gauge family covers exactly
the symmetric interior pairs, in agreement with \eqref{eq:hyperbolicbalance}.

At the nodes $(-a,a)$ the exact pairing has the matrix representation
\begin{equation}\label{eq:H2}
 H_a=\frac1{4a^2}
 \begin{pmatrix}
 1+a^2&3a^2-1\\
 3a^2-1&1+a^2
 \end{pmatrix}.
\end{equation}
It is diagonal if and only if $a^2=1/3$, equivalently $t=0$.  In that case
$H_a=I_2$ and the nodes are $\pm1/\sqrt3$.  For every $0<a<1$, the eigenvalues of $H_a$
are
\[
 1,\qquad \frac{1-a^2}{2a^2},
\]
so the matrix is positive definite even when it is nonlocal.  This example already separates
positive full-matrix power evaluation from positive scalar quadrature.

\subsection{A balanced asymmetric triple}

For $N=3$, the node condition is
\begin{equation}\label{eq:N3condition}
 x_1+x_2+x_3+x_1x_2x_3=0.
\end{equation}
Indeed, this is the expansion of \eqref{eq:productbalance}.  Given $x_1,x_2$ with
$1+x_1x_2\neq0$, one may set
\[
 x_3=-\frac{x_1+x_2}{1+x_1x_2}.
\]
The rational triple
\begin{equation}\label{eq:rationaltriple}
 \left(-\frac{19}{20},\frac25,\frac{55}{62}\right)
\end{equation}
is admissible and asymmetric.  Its ordinary interpolatory weights for Lebesgue integration
are
\[
 \frac{51200}{92259},\qquad
 \frac{18950}{12231},\qquad
 -\frac{53816}{515967}.
\]
The negative third weight shows that Stokes admissibility alone does not imply positivity of
the ordinary quadrature associated with the same roots.

\section{Nodal representation and the rigidity of locality}

\subsection{Every simple real kernel gives a full nodal matrix}

Suppose $\beta$ has full rank and its degree-$N$ kernel polynomial has distinct real roots
$X=(x_1,\ldots,x_N)$.  Define evaluation maps
\[
 V_Ep=(p(x_1),\ldots,p(x_N))^{\mathsf T},
 \qquad
 V_Fr=(r(x_1),\ldots,r(x_N))^{\mathsf T}.
\]
Then $V_F:F_N\to\R^N$ is an isomorphism, while
$\Ker V_E=\spanof\{\kappa_\beta\}=\Ker\beta^\flat$.

\begin{proposition}[Unique full nodal representation]\label{prop:fullnodal}
There exists a unique invertible matrix $H_\beta\in\R^{N\times N}$ such that
\begin{equation}\label{eq:fullnodal}
 \boxed{\beta(p,r)=(V_Ep)^{\mathsf T}H_\beta(V_Fr),}
 \qquad p\in E_N,\quad r\in F_N.
\end{equation}
\end{proposition}

\begin{proof}
The pairing descends to a nondegenerate bilinear form on
$(E_N/\spanof\{\kappa_\beta\})\times F_N$.  Both factors are identified with $\R^N$ by
$V_E$ and $V_F$, so a unique invertible representing matrix exists.
\end{proof}

The matrix $H_\beta$ couples different nodes.  A classical scalar quadrature corresponds to
the exceptional case in which it is diagonal.

\subsection{A polynomial spanning lemma}

\begin{lemma}\label{lem:derivativeproducts}
\begin{equation}\label{eq:spanproducts}
 \spanof\{(pq)':p,q\in\Pp_N\}=\Pp_{2N-1}.
\end{equation}
\end{lemma}

\begin{proof}
For $0\leq k\leq2N-1$, let $j=k+1$ and choose nonnegative integers $a,b\leq N$ with
$a+b=j$.  Then
\[
 x^k=\frac1{k+1}(x^ax^b)'.
\]
Thus every monomial of degree at most $2N-1$ belongs to the displayed span.
\end{proof}

\begin{theorem}[Scalar moment rigidity]\label{thm:momentrigidity}
Suppose a Stokes-compatible pairing has the product form
\begin{equation}\label{eq:momentform}
 \beta(p,r)=\mathcal L(pr)
\end{equation}
for a linear functional $\mathcal L$ on $\Pp_{2N-1}$.  Then
\begin{equation}\label{eq:Llebesgue}
 \mathcal L(f)=\int_{-1}^{1}f(x)\dd x,
 \qquad f\in\Pp_{2N-1},
\end{equation}
and hence $\beta=\beta_0$.
\end{theorem}

\begin{proof}
Stokes compatibility gives
\[
 \mathcal L((pq)')=\bracket{pq}=\int_{-1}^{1}(pq)'\dd x.
\]
Lemma~\ref{lem:derivativeproducts} implies equality of the two functionals on all of
$\Pp_{2N-1}$.  Products $pr$, with $p\in\Pp_N$ and $r\in\Pp_{N-1}$, belong to this space,
so the pairings coincide.
\end{proof}

This theorem already rules out a nontrivial weighted moment interpretation inside the fixed
unweighted Stokes class.

\subsection{Gauss--Legendre is the unique local scalar gauge}

\begin{theorem}[Rigidity of local scalar realization]\label{thm:localrigidity}
Let $\beta:E_N\times F_N\to\R$ be full rank and Stokes-compatible.  Suppose there exist
$N$ distinct real points $x_i$ and nonzero scalars $w_i$ such that
\begin{equation}\label{eq:localpairing}
 \beta(p,r)=\sum_{i=1}^Nw_i p(x_i)r(x_i),
 \qquad p\in E_N,\quad r\in F_N.
\end{equation}
Then:
\begin{enumerate}[label=\textup{(\roman*)}]
\item the rule $\sum_iw_if(x_i)$ is exact on $\Pp_{2N-1}$;
\item the node polynomial is proportional to the Legendre polynomial $P_N$;
\item the $x_i$ are the Gauss--Legendre nodes and the $w_i$ are the positive Gauss weights;
\item $\beta=\beta_0$, equivalently $\omega=0$.
\end{enumerate}
Conversely, $\beta_0$ has the representation \eqref{eq:localpairing} at the Gauss--Legendre
nodes and weights.
\end{theorem}

\begin{proof}
Define $\mathcal Q(f)=\sum_iw_if(x_i)$.  Applying Stokes to $p,q\in\Pp_N$ gives
\[
 \mathcal Q((pq)')=\bracket{pq}=\int_{-1}^{1}(pq)'\dd x.
\]
Lemma~\ref{lem:derivativeproducts} proves exactness on $\Pp_{2N-1}$.

Let $\pi_X(x)=\prod_i(x-x_i)$.  For every $r\in\Pp_{N-1}$,
\[
 \int_{-1}^{1}\pi_Xr\dd x
 =\sum_iw_i\pi_X(x_i)r(x_i)=0.
\]
Thus $\pi_X$ is the degree-$N$ polynomial orthogonal to $\Pp_{N-1}$ and is therefore
proportional to $P_N$.  The nodes are its roots.  The weights are uniquely determined by
exactness on $\Pp_{N-1}$.  If $\ell_i$ are the cardinal polynomials, exactness applied to
$\ell_i^2\in\Pp_{2N-2}$ gives
\[
 w_i=\int_{-1}^{1}\ell_i(x)^2\dd x>0.
\]
Finally, $pr\in\Pp_{2N-1}$, so \eqref{eq:localpairing} equals $\int pr$ for every admissible
pair.  This is $\beta_0$.  The converse is ordinary Gaussian exactness.
\end{proof}

\begin{corollary}\label{cor:diagonaliffzero}
For a full-rank pairing with a simple real degree-$N$ kernel, the matrix $H_\beta$ in
\eqref{eq:fullnodal} is diagonal if and only if $\omega=0$.  Hence every nonzero gauge is
nonlocal in its intrinsic kernel-node chart.
\end{corollary}

\begin{proof}
If $H_\beta$ is diagonal, its diagonal entries are nonzero because $H_\beta$ is invertible;
Theorem~\ref{thm:localrigidity} then gives $\beta=\beta_0$ and $\omega=0$.  Conversely,
the Gauss--Legendre formula represents $\beta_0$ by its positive diagonal weight matrix.
\end{proof}

The conclusion remains true when a nonzero gauge happens to retain $P_N$ as its kernel.
The nodes alone do not determine the pairing because the kernel map has positive-dimensional
fibres for $N\geq3$.

\section{Minimum-norm solution of a prescribed-kernel constraint}\label{sec:metric}

The affine Stokes structure alone carries no norm, so it cannot select a minimum gauge.
Here we use the distinguished $L^2$ reference to make such a selection.  Let
\[
 E_N^0=\left\{p\in E_N:\int_{-1}^{1}p(x)\dd x=0\right\}.
\]
The $L^2$ product identifies $E_N/\R$ with $E_N^0$ and induces a Euclidean norm
$\|\cdot\|_{\Lambda^2}$ on $\Lambda^2(E_N/\R)^*$.  Write
$Q_N^{(0)}=E_N/\R$ for this $L^2$ quotient and denote its inner product by
$\langle\cdot,\cdot\rangle_0$.  Every two-form $\omega$ has a unique skew-adjoint
representative $K_\omega\in\operatorname{End}(Q_N^{(0)})$,
\begin{equation}\label{eq:quotientriesz}
 \omega(p,q)=\langle[p],K_\omega[q]\rangle_0,
 \qquad
 \|\omega\|_{\Lambda^2}^2=\tfrac12\|K_\omega\|_{\mathrm{HS}}^2.
\end{equation}
Denote by $\Delta\beta_\omega$ the unique mixed pairing satisfying
\begin{equation}\label{eq:deltabetaomega}
 \Delta\beta_\omega(p,q')=\omega(p,q).
\end{equation}
For $\beta=\beta_0+\Delta\beta_{\omega_\beta}$, we use the induced gauge norm
\begin{equation}\label{eq:pairinggaugenorm}
 \boxed{
 \|\beta-\beta_0\|:=\|\omega_\beta\|_{\Lambda^2}
 =\frac1{\sqrt2}\|K_{\omega_\beta}\|_{\mathrm{HS}}.}
\end{equation}

For an admissible nonconstant polynomial $\kappa$, define
$u=[\kappa]\in Q_N^{(0)}$ and let $g\in Q_N^{(0)}$ be the Riesz representative of the
functional
\begin{equation}\label{eq:gdefinition}
 \langle g,[q]\rangle_0=\int_{-1}^{1}\kappa(x)q'(x)\dd x.
\end{equation}
The endpoint condition gives
\begin{equation}\label{eq:guorthogonal}
 \langle g,u\rangle_0
 =\int_{-1}^{1}\kappa\kappa'\dd x
 =\tfrac12\bigl(\kappa(1)^2-\kappa(-1)^2\bigr)=0.
\end{equation}
For $v\in Q_N^{(0)}$, the notation $v^\flat$ denotes the Riesz covector
$w\mapsto\langle v,w\rangle_0$.

Let $P_N$ denote the Legendre polynomial normalized by $P_N(1)=1$, and set
\begin{equation}\label{eq:wdefinition}
 s=\kappa(1),\qquad w=sP_N-\kappa.
\end{equation}
Then $w(1)=w(-1)=0$.

\begin{lemma}[Explicit Riesz vector]\label{lem:explicitg}
The vector in \eqref{eq:gdefinition} is
\begin{equation}\label{eq:gexplicit}
 \boxed{g=w'=(sP_N-\kappa)'.}
\end{equation}
\end{lemma}

\begin{proof}
For every $q\in E_N$, orthogonality of $P_N$ to $q'\in F_N$ gives
\begin{equation}\label{eq:endpointlegendreidentity}
 q(1)-(-1)^Nq(-1)=\int_{-1}^{1}P_N'(x)q(x)\dd x.
\end{equation}
Integration by parts and
$\kappa(-1)=(-1)^Ns$ therefore yield
\[
 \int_{-1}^{1}\kappa q'
 =s\bigl(q(1)-(-1)^Nq(-1)\bigr)-\int_{-1}^{1}\kappa'q
 =\int_{-1}^{1}w'q.
\]
Moreover, $\int_{-1}^{1}w'=w(1)-w(-1)=0$, so $w'\in E_N^0$ is precisely the Riesz
representative in the zero-mean realization of $Q_N^{(0)}$.
\end{proof}

Let
\begin{equation}\label{eq:gaugefiberclosure}
 \mathcal V_\kappa=
 \left\{\omega\in\Lambda^2(E_N/\R)^*:
  (\beta_0+\Delta\beta_\omega)(\kappa,r)=0\ \text{for every }r\in F_N
 \right\}.
\end{equation}
The full-rank locus in $\mathcal V_\kappa$ is defined by the nonvanishing of a maximal minor,
so it is a relative Zariski-open subset.  It is nonempty by
Theorem~\ref{thm:kernelimage}; since $\mathcal V_\kappa$ is an affine space and hence
irreducible, that locus is dense.  Thus $\mathcal V_\kappa$ is exactly the affine closure of
the full-rank fibre with prescribed kernel line.

\begin{theorem}[Closed formula for the minimum gauge]\label{thm:minimumgauge}
The radical constraint $\omega\in\mathcal V_\kappa$ is equivalent to
\begin{equation}\label{eq:KuEqualsg}
 K_\omega u=g.
\end{equation}
Its unique minimum-norm solution is
\begin{equation}\label{eq:Kcanonical}
 \boxed{
 K_{\mathrm{can}}
 =\frac{g\otimes u^\flat-u\otimes g^\flat}{\|u\|_0^2},
 \qquad
 \|\omega_{\mathrm{can}}\|_{\Lambda^2}
 =\frac{\|g\|_0}{\|u\|_0}.}
\end{equation}
Every solution is uniquely of the form
\begin{equation}\label{eq:allfibreoperators}
 K_\omega=K_{\mathrm{can}}+L,
 \qquad L^*=-L,\quad Lu=0,
\end{equation}
and the two terms are Hilbert--Schmidt orthogonal.  In particular, the residual affine
freedom has dimension $(N-1)(N-2)/2$.  The construction depends only on the projective line
$[\kappa]$ and on the chosen $L^2$ quotient metric.
\end{theorem}

\begin{proof}
For $r=q'$, the radical equation is
\[
 0=\int\kappa q'+\omega(\kappa,q)
  =\langle g,[q]\rangle_0+\langle u,K_\omega[q]\rangle_0
  =\langle g-K_\omega u,[q]\rangle_0,
\]
which proves \eqref{eq:KuEqualsg}.  Orthogonality \eqref{eq:guorthogonal} shows that the
operator in \eqref{eq:Kcanonical} is skew and sends $u$ to $g$.  The difference of any two
solutions is a skew operator annihilating $u$, proving \eqref{eq:allfibreoperators}.  A direct
trace computation gives
$\langle K_{\mathrm{can}},L\rangle_{\mathrm{HS}}=0$ whenever $Lu=0$; hence Pythagoras proves
minimality and uniqueness.  The operator acts only on $\spanof\{u,g\}$, so
$\tfrac12\|K_{\mathrm{can}}\|_{\mathrm{HS}}^2=\|g\|_0^2/\|u\|_0^2$.
Finally, $L$ is an arbitrary skew operator on $u^\perp$, which gives the dimension.  Both
$u$ and $g$ scale with $\kappa$, so \eqref{eq:Kcanonical} is projectively invariant.
\end{proof}

Let $\omega_{\mathrm{can}}(\kappa)$ be the two-form represented by
$K_{\mathrm{can}}$ and put
\begin{equation}\label{eq:betacankappa}
 \beta_{\mathrm{can}}(\kappa)
 =\beta_0+\Delta\beta_{\omega_{\mathrm{can}}(\kappa)}.
\end{equation}

\begin{theorem}[Global rank stratification of the canonical representative]
\label{thm:globalcanonicalrank}
For every admissible nonconstant $\kappa\in E_N$,
\begin{equation}\label{eq:globalcanonicalrank}
 \boxed{
 \rank\beta_{\mathrm{can}}(\kappa)
 =
 \begin{cases}
  N,   & \deg\kappa=N,\\
  N-2, & \deg\kappa<N.
 \end{cases}}
\end{equation}
Put
\begin{equation}\label{eq:acanoperator}
 a_{\mathrm{can}}(p,q)=\beta_{\mathrm{can}}(\kappa)(p,q'),
 \qquad
 T_{\mathrm{can}}q=\iota(q')+K_{\mathrm{can}}[q],
\end{equation}
where $\iota:F_N\hookrightarrow E_N$ is the natural inclusion.
If $\deg\kappa=N$, then
\begin{equation}\label{eq:canonicalradicalexact}
 \Rad_Ra_{\mathrm{can}}=\spanof\{1\},
 \qquad
 \Rad_La_{\mathrm{can}}=\Rad_L\beta_{\mathrm{can}}(\kappa)
 =\spanof\{\kappa\}.
\end{equation}
If $\deg\kappa<N$, define the zero-mean representative and its centred primitive by
\begin{equation}\label{eq:kappa0eta}
 \bar\kappa=\frac12\int_{-1}^{1}\kappa(x)\dd x,
 \qquad
 \kappa_0=\kappa-\bar\kappa,
 \qquad
 \eta(x)=\int_{-1}^{x}\kappa_0(t)\dd t.
\end{equation}
Then, with $w=sP_N-\kappa$ as in \eqref{eq:wdefinition},
\begin{equation}\label{eq:lowerdegreeradicals}
 \boxed{
 \Rad_Ra_{\mathrm{can}}=\spanof\{1,w,\eta\},
 \qquad
 \Rad_La_{\mathrm{can}}=\spanof\{\kappa,w,\eta\}.}
\end{equation}
Thus the lower-degree strata have an exact rank defect two.  They occur only for $N\geq3$.
\end{theorem}

\begin{proof}
Identifying quotient classes with their zero-mean representatives gives
\[
 a_{\mathrm{can}}(p,q)=\langle p,T_{\mathrm{can}}q\rangle_{L^2}.
\]
The $L^2$ form on $E_N$ is nondegenerate, and $d:E_N\to F_N$ is onto; hence
\begin{equation}\label{eq:rankbetaT}
 \rank\beta_{\mathrm{can}}(\kappa)=\rank a_{\mathrm{can}}
 =\rank T_{\mathrm{can}}.
\end{equation}

Assume first that $\deg\kappa<N$.  Orthogonality to $P_N$ gives
\begin{equation}\label{eq:lowerdegreeinnerproducts}
 \langle u,[w]\rangle_0=-\|u\|_0^2,
 \qquad
 \langle g,[w]\rangle_0=\int_{-1}^{1}w'w\dd x=0.
\end{equation}
Formula \eqref{eq:Kcanonical} therefore yields $K_{\mathrm{can}}[w]=-g$, and thus
$T_{\mathrm{can}}w=0$.  Since $\kappa_0$ has zero mean, the polynomial $\eta$ in
\eqref{eq:kappa0eta} vanishes at both endpoints.  Moreover,
\begin{equation}\label{eq:etainnerproducts}
 \langle u,[\eta]\rangle_0
 =\int_{-1}^{1}\eta'\eta\dd x=0,
 \qquad
 \langle g,[\eta]\rangle_0
 =\int_{-1}^{1}w'\eta\dd x
 =-\int_{-1}^{1}w\kappa_0\dd x
 =\|u\|_0^2.
\end{equation}
Consequently $K_{\mathrm{can}}[\eta]=-u$ and
$T_{\mathrm{can}}\eta=\eta'-\kappa_0=0$.  Also $T_{\mathrm{can}}1=0$.

The polynomials $1,w,\eta$ are linearly independent.  Indeed, evaluating a possible
relation at either endpoint first removes its constant term.  If $\deg\kappa\leq N-2$,
then $w$ and $\eta$ have different degrees.  If $\deg\kappa=N-1$, proportionality of
$w$ and $\eta$ would imply proportionality of the nonzero vectors $g=w'$ and $u=\kappa_0$,
contrary to $\langle g,u\rangle_0=0$.  Hence
\begin{equation}\label{eq:rankTupper}
 \rank T_{\mathrm{can}}\leq N-2.
\end{equation}
On the other hand, $\iota d:E_N\to E_N$ has rank $N$, while $K_{\mathrm{can}}$ has rank two
because $u$ and $g$ are nonzero and orthogonal.  The rank inequality gives
\begin{equation}\label{eq:rankTlower}
 \rank T_{\mathrm{can}}=\rank(\iota d+K_{\mathrm{can}})
 \geq\rank(\iota d)-\rank K_{\mathrm{can}}=N-2.
\end{equation}
Equations \eqref{eq:rankbetaT}--\eqref{eq:rankTlower} prove the lower branch of
\eqref{eq:globalcanonicalrank} and the right-radical identity in
\eqref{eq:lowerdegreeradicals}.

Both $w$ and $\eta$ vanish at the endpoints.  Since they lie in the right radical, the
Stokes identity puts them in the left radical as well.  The polynomial $\kappa$ is in the
left radical by construction and is independent of $w,\eta$ because its endpoint value is
nonzero.  The left nullity is three, so these vectors give the second identity in
\eqref{eq:lowerdegreeradicals}.  For $N=1$ there is no nonconstant lower-degree polynomial;
for $N=2$, the even endpoint condition forces every affine admissible polynomial to be
constant.  Hence the lower-degree branch can occur only for $N\geq3$.

Assume now that $\deg\kappa=N$ and let $p$ belong to the left radical.  Since both $p$ and
$\kappa$ annihilate $F_N$,
\[
 0=b(\kappa,p)
 =s\bigl(p(1)-(-1)^Np(-1)\bigr).
\]
Put $t=p(1)$ and $h=tP_N-p$.  Then $h(1)=h(-1)=0$, and the calculation of
Lemma~\ref{lem:explicitg}, now applied to $p$, shows that its Riesz vector is $h'$.  The
radical equation is therefore
\begin{equation}\label{eq:pradicalcanonical}
 h'=K_{\mathrm{can}}[p]
 =\frac{g\langle u,[p]\rangle_0-u\langle g,[p]\rangle_0}{\|u\|_0^2}.
\end{equation}
The polynomials $h'$ and $g$ have degree at most $N-1$, whereas the zero-mean representative
of $u$ has degree exactly $N$.  Comparing the degree-$N$ component in
\eqref{eq:pradicalcanonical} gives
$\langle g,[p]\rangle_0=0$.  With
\[
 a=\frac{\langle u,[p]\rangle_0}{\|u\|_0^2},
\]
we obtain $h'=ag$.  The polynomials $h$ and $aw$ have the same derivative and both vanish at
the endpoints; hence $h=aw$.  Consequently,
\begin{equation}\label{eq:pdecompositioncanonical}
 p=a\kappa+(t-as)P_N.
\end{equation}
Taking the inner product with $u$ and using the definition of $a$ gives
\[
 (t-as)\langle u,[P_N]\rangle_0=0.
\]
Because $\kappa$ has exact degree $N$, its coefficient along $P_N$ is nonzero, so
$\langle u,[P_N]\rangle_0\neq0$.  Thus $t=as$ and $p=a\kappa$.  The left radical is exactly
one dimensional, which is equivalent to full rank.  The right radical of
$a_{\mathrm{can}}$ contains the constants and, by rank--nullity, is exactly
$\spanof\{1\}$.
\end{proof}

\begin{corollary}[Attainment of the minimum gauge in a full-rank fibre]
\label{cor:fullrankminimum}
Let $\mathcal F_\kappa^{\mathrm{fr}}$ be the full-rank fibre over an admissible kernel line
$[\kappa]$.  Then
\begin{equation}\label{eq:fullrankminimum}
 \inf_{\beta\in\mathcal F_\kappa^{\mathrm{fr}}}
 \|\beta-\beta_0\|
 =\|\beta_{\mathrm{can}}(\kappa)-\beta_0\|.
\end{equation}
The infimum is attained if and only if $\deg\kappa=N$; in that case its unique minimizer is
$\beta_{\mathrm{can}}(\kappa)$.  If $\deg\kappa<N$, the infimum is not attained.
\end{corollary}

\begin{proof}
The prescribed-kernel constraint is affine and closed, and its unique minimum-norm element
is $\beta_{\mathrm{can}}(\kappa)$.  Full rank is an open dense condition in every nonempty
fibre, so the infimum over the full-rank part equals the minimum over its closure.  By
Theorem~\ref{thm:globalcanonicalrank}, the canonical representative is full rank exactly
when $\deg\kappa=N$.  Strict convexity of the squared norm gives uniqueness in that case;
in lower degree the unique minimizer lies outside the full-rank locus, so the common
infimum cannot be attained there.
\end{proof}

\begin{corollary}[Rectangular right radical]\label{cor:rectangularradical}
For a mixed pairing write
\[
 \Rad_R\beta=\{r\in F_N:\beta(p,r)=0\ \text{for every }p\in E_N\}.
\]
The canonical representative satisfies
\begin{equation}\label{eq:rectangularradical}
 \boxed{
 \Rad_R\beta_{\mathrm{can}}(\kappa)=
 \begin{cases}
  \{0\},&\deg\kappa=N,\\
  \spanof\{g,\kappa_0\},&\deg\kappa<N,
 \end{cases}}
\end{equation}
where $g=w'$ and $\kappa_0=\eta'$ are nonzero and $L^2$-orthogonal in the second case.
Thus the rectangular left and right nullities are respectively $(1,0)$ in exact degree and
$(3,2)$ in lower degree.
\end{corollary}

\begin{proof}
Surjectivity of $d:E_N\to F_N$ identifies $F_N$ with $E_N/\R$.  The right radical of
$\beta_{\mathrm{can}}$ is therefore the derivative of the right radical of
$a_{\mathrm{can}}$.  Equations \eqref{eq:canonicalradicalexact} and
\eqref{eq:lowerdegreeradicals} give respectively $\{0\}$ and
$\spanof\{w',\eta'\}$.  Nonvanishing and orthogonality follow from the proof of
Theorem~\ref{thm:globalcanonicalrank}.
\end{proof}

\begin{lemma}[Rectangular singular-value opening]\label{lem:rectangularopening}
Let $M(t):X\to Y$ be a smooth family between finite-dimensional Hilbert spaces with
$\dim Y\geq\dim X=n$, and suppose $\rank M(0)=n-q$.  Put
\begin{equation}\label{eq:rectangularcrossing}
 R=\Ker M(0),\qquad L=\Ker M(0)^*,\qquad
 G=\Pi_LM'(0)|_R:R\longrightarrow L.
\end{equation}
If $G$ is injective, then there are $c,C,\delta>0$ such that, for $0<|t|<\delta$,
\begin{equation}\label{eq:rectangularsvbounds}
 c|t|\leq\sigma_j(M(t))\leq C|t|,
 \qquad j=n-q+1,\ldots,n,
\end{equation}
where the singular values are the eigenvalues of $(M(t)^*M(t))^{1/2}$, counted with
multiplicity and ordered decreasingly.  The other $n-q$ singular values remain uniformly
separated from zero.
\end{lemma}

\begin{proof}
Use the orthogonal decompositions $X=R^\perp\oplus R$ and
$Y=\im M(0)\oplus L$.  In the corresponding blocks,
\[
 M(t)=\begin{pmatrix}A(t)&B(t)\\ C(t)&D(t)\end{pmatrix},
 \qquad A(0):R^\perp\longrightarrow\im M(0)
 \ \text{is an isomorphism},
\]
and $B(0)=C(0)=D(0)=0$.  For small $t$, bounded invertible block eliminations reduce
$M(t)$ to $\diag(A(t),S(t))$, where
\[
 S(t)=D(t)-C(t)A(t)^{-1}B(t)=tG+O(t^2).
\]
Injectivity of $G$ gives two-sided linear singular-value bounds for $S(t)$.  These bounds
survive the eliminations because, for invertible $P,Q$ and every admissible index $j$,
\begin{equation}\label{eq:svelimination}
 \frac{\sigma_j(B)}{\|P^{-1}\|\,\|Q^{-1}\|}
 \leq\sigma_j(PBQ)
 \leq\|P\|\,\|Q\|\,\sigma_j(B).
\end{equation}
The singular values belonging to $A(t)$ remain separated from zero by continuity.
\end{proof}

\begin{theorem}[Transverse opening of the lower-degree defect]
\label{thm:transverseopening}
Let $\kappa_\star$ be an admissible lower-degree kernel, set
\begin{equation}\label{eq:transversedata}
 s=\kappa_\star(1),\quad
 u_\star=\kappa_\star-\frac12\int_{-1}^{1}\kappa_\star(x)\dd x,\quad
 A=\|u_\star\|_{L^2}^2,\quad B=\|P_N\|_{L^2}^2,
\end{equation}
and define $\kappa_t=\kappa_\star+tP_N$.  For $t\neq-s$, this remains endpoint-admissible.
Let $a_t$ be the composed canonical form associated with $\kappa_t$, and set
$T_tq=\iota(q')+K_{\mathrm{can}}(\kappa_t)[q]$.  Let
\begin{equation}\label{eq:transversewetag}
 w=sP_N-\kappa_\star,\qquad g=w',\qquad
 \eta(x)=\int_{-1}^{x}u_\star(\tau)\dd\tau.
\end{equation}
Then the restriction of $a_t$ to the ordered basis $(w,\eta)$ is exactly
\begin{equation}\label{eq:transverseblock}
 \boxed{
 [a_t]_{(w,\eta)}
 =\frac{ABt(s+t)}{A+Bt^2}
 \begin{pmatrix}0&-1\\1&0\end{pmatrix}.}
\end{equation}
In particular,
\begin{equation}\label{eq:transversederivative}
 \left.\frac{\dd}{\dd t}[a_t]_{(w,\eta)}\right|_{t=0}
 =sB\begin{pmatrix}0&-1\\1&0\end{pmatrix},
\end{equation}
so the rank-two defect opens transversely.  With respect to any fixed Hilbert structures,
equivalently in fixed orthonormal bases, the two vanishing singular values
$\sigma_{N-1}(\beta_t)$ and $\sigma_N(\beta_t)$ of the rectangular pairing
satisfy
\begin{equation}\label{eq:transversesvbounds}
 c|t|\leq\sigma_j(\beta_t)\leq C|t|,
 \qquad j=N-1,N,
\end{equation}
for some $c,C>0$ and all sufficiently small $t\neq0$, while the remaining singular values
stay bounded away from zero.  Consequently
$\operatorname{cond}(\beta_t)=\sigma_1(\beta_t)/\sigma_N(\beta_t)$ is of order $|t|^{-1}$.
\end{theorem}

\begin{proof}
Since $u_\star$, $g$, and $w'$ have degree below $N$, Legendre orthogonality gives
\begin{equation}\label{eq:transverseorthogonality}
 \langle u_\star,P_N\rangle=\langle g,P_N\rangle=0,
 \qquad \|u_\star+tP_N\|^2=A+Bt^2.
\end{equation}
Moreover,
\begin{equation}\label{eq:transverseinnerproducts}
 \langle u_\star,w\rangle=-A,
 \qquad \langle P_N,w\rangle=sB,
 \qquad \langle g,w\rangle=0,
 \qquad \langle g,\eta\rangle=A.
\end{equation}
The Legendre defect associated with $\kappa_t$ is still $w$, hence its Riesz vector is still
$g$.  Substitution into \eqref{eq:Kcanonical} yields
\begin{equation}\label{eq:Ttw}
 T_tw=\frac{Bt(s+t)}{A+Bt^2}\,g.
\end{equation}
Because $w$ and $\eta$ vanish at both endpoints, the restriction of $a_t$ to their span is
skew.  Equations \eqref{eq:transverseinnerproducts}--\eqref{eq:Ttw} give
$a_t(\eta,w)=ABt(s+t)/(A+Bt^2)$, proving \eqref{eq:transverseblock} and
\eqref{eq:transversederivative}.  At $t=0$, the corresponding left null vectors are
$w,\eta$ and the right null vectors of the mixed pairing are $g,\eta'=u_\star$ by
Corollary~\ref{cor:rectangularradical}.  The derivative block is nondegenerate because
$sB\neq0$.  Therefore the crossing map in Lemma~\ref{lem:rectangularopening} is injective
on the two-dimensional right kernel.  The lemma gives \eqref{eq:transversesvbounds} and the
claimed condition-number estimate.
\end{proof}

\begin{remark}[Pairing section versus nodal section]
The projectively invariant map $[\kappa]\mapsto\beta_{\mathrm{can}}(\kappa)$ is a global
canonical pairing section on the entire admissible degree-$N$ kernel stratum, including
classes with multiple or nonreal roots.  A nodal matrix requires $N$ distinct roots, and an
ordered interior node configuration requires those roots to be real and simple.  Thus the
canonical interior nodal section below is the restriction of the pairing section to the
simple interior real-root stratum.
\end{remark}

At the Legendre line, $w=0$ and $g=0$, so the canonical gauge is the zero gauge.

For an admissible ordered interior node set $X=(x_1,\ldots,x_N)$, put
$\pi_X(x)=\prod_i(x-x_i)$, let $\omega_{\mathrm{can}}(X)$ be the two-form represented by
\eqref{eq:Kcanonical} with $\kappa=\pi_X$, and define
$\beta_{\mathrm{can}}(X)=\beta_0+\Delta\beta_{\omega_{\mathrm{can}}(X)}$.

\begin{corollary}[Global canonical nodal section and local hyperbolic stability]
\label{thm:localstability}
For every simple admissible interior configuration $X$, the pairing
$\beta_{\mathrm{can}}(X)$ has full rank and its intrinsic nodal matrix
$H_{\mathrm{can}}(X)$ is well defined.  Let
$X_G=(\xi_1,\ldots,\xi_N)$ be the ordered Gauss--Legendre configuration and put
$y_i=\operatorname{artanh}x_i$, $y_i^G=\operatorname{artanh}\xi_i$.  There exist a
neighbourhood $U$ of $X_G$ in the admissible configuration manifold and a constant $C_U$
such that, for every $X\in U$,
\begin{equation}\label{eq:localstabilitybound}
 \|\omega_{\mathrm{can}}(X)\|_{\Lambda^2}
 +\|\operatorname{off}H_{\mathrm{can}}(X)\|_{\mathrm F}
 \leq C_U\|y-y^G\|_2.
\end{equation}
Here $\operatorname{off}H=H-\diag(H_{11},\ldots,H_{NN})$.
\end{corollary}

\begin{proof}
In fixed polynomial coordinates, the coefficients of $\pi_X$, and hence the vectors $u(X)$
and $g(X)$ in \eqref{eq:gdefinition}, depend smoothly on $X$.  Since $u(X)\neq0$, formula
\eqref{eq:Kcanonical} makes $\omega_{\mathrm{can}}(X)$ smooth on the degree-$N$ admissible
stratum.  At $X_G$, the Legendre polynomial is orthogonal to $F_N$, so $g(X_G)=0$ and
$\omega_{\mathrm{can}}(X_G)=0$.

Theorem~\ref{thm:globalcanonicalrank} gives full rank because $\deg\pi_X=N$.  On the
simple-node stratum, the evaluation isomorphisms in Proposition~\ref{prop:fullnodal} depend
smoothly on $X$; therefore $H_{\mathrm{can}}(X)$ is smooth.  At $X_G$ it is the diagonal matrix of
positive Gauss weights.  Thus both maps on the left of \eqref{eq:localstabilitybound} vanish
at $X_G$ and are locally Lipschitz.  Finally $x\mapsto\operatorname{artanh}x$ is a smooth
local diffeomorphism, so rapidity and Euclidean coordinates are locally bi-Lipschitz.
\end{proof}

The estimate cannot be extended uniformly to the whole admissible chamber.

\begin{proposition}[Failure of a global hyperbolic bound]\label{prop:noglobalbound}
There is no constant $C$ such that
\begin{equation}\label{eq:falseglobalbound}
 \|\operatorname{off}H_{\mathrm{can}}(X)\|_{\mathrm F}
 \leq C\|\operatorname{artanh}X-\operatorname{artanh}X_G\|_2
\end{equation}
for every simple admissible interior configuration, even when $N=2$.
\end{proposition}

\begin{proof}
For $N=2$ every fibre has dimension zero, so the pairing at $X(a)=(-a,a)$ is forced and its
nodal matrix is \eqref{eq:H2}.  Hence
\[
 \|\operatorname{off}H_a\|_{\mathrm F}
 =\frac{\sqrt2\,|3a^2-1|}{4a^2}\longrightarrow\infty
 \qquad(a\downarrow0).
\]
On the other hand, with $c=\operatorname{artanh}(1/\sqrt3)$,
\[
 \|\operatorname{artanh}X(a)-\operatorname{artanh}X_G\|_2
 =\sqrt2\,|\operatorname{artanh}a-c|\longrightarrow\sqrt2\,c<\infty.
\]
Thus \eqref{eq:falseglobalbound} is impossible.
\end{proof}

The degeneration $a\downarrow0$ is a node collision at finite hyperbolic distance.  Any
global stability estimate must therefore contain an additional factor measuring separation,
the nodal discriminant, or the condition number of the evaluation map.  Hyperbolic distance
controls displacement from Gauss but cannot by itself control collision geometry.

\section{Ordinary quadrature attached only to the roots}

Every set of distinct real nodes admits an interpolatory rule for Lebesgue integration,
independently of the pairing gauge.  Let $\ell_i$ be the Lagrange basis and define
\begin{equation}\label{eq:interpweights}
 \widehat w_i=\int_{-1}^{1}\ell_i(x)\dd x,
 \qquad
 Q_X(f)=\sum_i\widehat w_i f(x_i).
\end{equation}
The rule is exact on $\Pp_{N-1}$, but it need not represent $\beta$ and its weights need not
be positive.

\begin{proposition}[Degree of the interpolatory rule]\label{prop:degree}
Let $\kappa_X=\prod_i(x-x_i)$ and let $s\in\{0,\ldots,N\}$ be maximal such that
\begin{equation}\label{eq:orthocriterion}
 \int_{-1}^{1}\kappa_X(x)q(x)\dd x=0,
 \qquad q\in\Pp_{s-1},
\end{equation}
where the condition is empty for $s=0$.  Then $Q_X$ has degree of exactness
$N+s-1$.  If $\kappa_X$ is Stokes-admissible but not proportional to $P_N$, its degree is at
most $2N-3$.
\end{proposition}

\begin{proof}
For $f\in\Pp_{N+s-1}$, polynomial division gives
$f=\kappa_Xq+r$ with $q\in\Pp_{s-1}$ and $r\in\Pp_{N-1}$.  Since $Q_X(\kappa_Xq)=0$ and
$Q_X(r)=\int r$, exactness is equivalent to \eqref{eq:orthocriterion}.  Maximality gives the
degree when $s<N$.  If $s=N$, the rule is exact through degree $2N-1$, while
$Q_X(\kappa_X^2)=0\neq\int_{-1}^{1}\kappa_X^2$ shows that it is not exact in degree $2N$.

If the degree were at least $2N-2$, then $\kappa_X\perp\Pp_{N-2}$, so
$\kappa_X=aP_N+bP_{N-1}$.  With the normalization $P_j(1)=1$, Stokes admissibility gives
\[
 a+b=(-1)^N\bigl(a(-1)^N+b(-1)^{N-1}\bigr)=a-b,
\]
hence $b=0$.
\end{proof}

Thus there are three distinct statements:
\begin{enumerate}
\item the roots of $\kappa_\beta$ are intrinsic to the gauged pairing;
\item those roots always support a full matrix representation of $\beta$;
\item they support a scalar quadrature representation of $\beta$ only in the zero gauge.
\end{enumerate}

\section{Orthogonal-polynomial families and the endpoint obstruction}

Every parity-symmetric polynomial family satisfies the endpoint condition degree by degree.
Consequently, Legendre, Chebyshev, Gegenbauer, symmetric Jacobi, and orthogonal polynomials
for an even positive measure can each be realized as kernel polynomials of suitable gauges,
provided their endpoint values are nonzero.  This statement concerns the kernel line only; it
does not identify the gauged pairing with the positive measure that originally defined the
family.

The endpoint constraint becomes especially informative for a coherent orthogonal hierarchy.

\begin{theorem}[Centered Favard recurrence]\label{thm:favard}
Let $(\kappa_n)_{n\geq0}$ be monic orthogonal polynomials for a positive measure supported on
$[-1,1]$, with recurrence
\begin{equation}\label{eq:favard}
 \kappa_{n+1}(x)=(x-a_n)\kappa_n(x)-b_n\kappa_{n-1}(x),
 \qquad b_n>0.
\end{equation}
Assume for every $n$ that
\begin{equation}\label{eq:hierarchyendpoint}
 \kappa_n(1)=(-1)^n\kappa_n(-1)\neq0.
\end{equation}
Then
\begin{equation}\label{eq:an0}
 a_n=0\qquad\text{for every }n.
\end{equation}
Hence the Jacobi matrix is centred and the orthogonality measure is symmetric.
\end{theorem}

\begin{proof}
Write $s_n=\kappa_n(1)$.  Evaluating \eqref{eq:favard} at $1$ gives
\[
 s_{n+1}=(1-a_n)s_n-b_ns_{n-1}.
\]
Evaluating at $-1$ and using \eqref{eq:hierarchyendpoint} gives
\[
 s_{n+1}=(1+a_n)s_n-b_ns_{n-1}.
\]
Since $s_n\neq0$, comparison yields $a_n=0$.  A compactly supported positive measure is
determined by its Jacobi coefficients; vanishing diagonal coefficients give invariance under
$x\mapsto-x$.
\end{proof}

The theorem explains why the fixed unweighted Stokes boundary is naturally compatible with
centred classical hierarchies.  General Jacobi families with unequal parameters require a
weighted or otherwise modified Stokes geometry.

\section{Relation with SBP and open Dirac structures}

\subsection{The skew freedom in SBP}

The algebraic SBP identity is
\begin{equation}\label{eq:sbp}
 Q+Q^{\mathsf T}=B.
\end{equation}
Its general solution is $Q=\tfrac12B+S$ with $S^{\mathsf T}=-S$.  In the rectangular
construction, $Q$ is replaced by
\[
 A=MD,
\]
and the derivative kernel additionally imposes $A(\cdot,1)=0$.  Theorem~\ref{thm:affine}
is therefore the exact quotient-aware version of the familiar skew freedom in SBP.

The containment of square SBP in the rectangular construction can be stated without a
choice of polynomial basis.

\begin{proposition}[Symmetric rectangularization of SBP]\label{thm:sbpinstance}
Let $\mathsf D\in\R^{n\times n}$ satisfy
\begin{equation}\label{eq:squaresbp}
 H\mathsf D+\mathsf D^{\mathsf T}H=B,
 \qquad H=H^{\mathsf T},
\end{equation}
and assume $\Ker\mathsf D=\spanof\{\bm 1\}$.  Put
$\mathsf F=\im\mathsf D$, let $J:\mathsf F\hookrightarrow\R^n$ be the inclusion, and let
$D:\R^n\to\mathsf F$ be the corestriction of $\mathsf D$, so that $\mathsf D=JD$.
Then
\begin{equation}\label{eq:rectfromsquare}
 M:=HJ:\mathsf F\longrightarrow(\R^n)^*
\end{equation}
satisfies the rectangular Stokes identity
\begin{equation}\label{eq:rectsbp}
 MD+D^{\mathsf T}M^{\mathsf T}=B.
\end{equation}
If $H$ is nonsingular, then $M$ has full column rank.  Conversely, a rectangular factorization
$\mathsf D=JD$, $M=HJ$ with symmetric $H$ turns \eqref{eq:rectsbp} into the square SBP
identity \eqref{eq:squaresbp}.
\end{proposition}

\begin{proof}
Substitution of $\mathsf D=JD$ and $M=HJ$ gives
\[
 MD+D^{\mathsf T}M^{\mathsf T}
 =HJD+D^{\mathsf T}J^{\mathsf T}H
 =H\mathsf D+\mathsf D^{\mathsf T}H=B.
\]
The inclusion $J$ has full column rank, so $HJ$ does as well when $H$ is nonsingular.  The
same identities read in reverse prove the converse.
\end{proof}

For nodal polynomial SBP, $\mathsf F$ is the sampled derivative space and $D$ is the exact
rectangular derivative after choosing coordinates in that space.  For FSBP, the same
argument applies with the sampled derivative image of the chosen function space, whenever
constants form the derivative kernel.  Proposition~\ref{thm:sbpinstance} is therefore an
algebraic containment statement: such a square identity is a symmetric factorization of the
mixed Stokes relation, with the additional restriction $M=HJ$.  It is a metric factorization
only when $H$ is positive definite.

The phrase \emph{dual-pairing SBP} is used elsewhere for paired forward and backward
difference operators acting on the same discrete grid \cite{williamsduru2024}.  That notion
should not be identified terminologically with the present mixed pairing
$E\times F\to\R$, whose two factors have different dimensions and are connected by the
surjective derivative.  The frameworks share a power-balance identity but organize the
discrete variables differently.

In the usual unreduced nodal coordinates, an SBP operator is written with a symmetric
positive norm matrix $P$:
\[
 PD_h+D_h^{\mathsf T}P=R^{\mathsf T}\Sigma R.
\]
Generalized SBP links positive norm matrices to quadrature
\cite{hickenzingg2013,delrey2014}; FSBP carries the same principle to arbitrary function
spaces \cite{glaubitz2023}.  The present classification concerns the prior rectangular level,
where no symmetric positive norm has yet been selected.

\subsection{Boundary completion}

Let $R:E_N\to\R^s$ be a collection of boundary traces.  The augmented power map
\[
 p\longmapsto(\beta^\flat p,Rp)
\]
is injective if and only if
\begin{equation}\label{eq:boundarycompletion}
 \Ker\beta^\flat\cap\Ker R=\{0\}.
\end{equation}
For a full-rank pairing this is equivalent to $R\kappa_\beta\neq0$.  The endpoint condition
\eqref{eq:endpointparity} ensures that either endpoint trace completes the kernel direction on
the degree-regular full-rank stratum.  This proves the injective boundary-completion mechanism
used in the 2012 construction.  It does not, by itself, prove maximal isotropy of a complete
open Dirac relation: that statement also requires specifying the full flow--effort relation,
its boundary-port pairing, and the ambient symmetric form.

Gauge-dependent kernel nodes and arbitrary collocation nodes must not be conflated.  Any
unisolvent grid may represent the flow polynomial.  The zeros of $\kappa_\beta$ instead
describe a quotient chart intrinsic to the chosen pairing, and Corollary~\ref{cor:diagonaliffzero}
shows that this chart is generally matrix-weighted rather than diagonal.

\section{Infinite-dimensional interpretation}\label{sec:continuum}

\subsection{The continuous affine class}

The finite-dimensional classification has an exact functional analogue.  Let
\[
 V=H^1(-1,1),\qquad W=L^2(-1,1),\qquad d=\partial_x.
\]
The derivative induces a Banach-space isomorphism
\[
 \overline d:V/\R\longrightarrow W.
\]

\begin{theorem}[Continuous Stokes pairing gauge]\label{thm:continuous}
The continuous bilinear pairings $\beta:V\times W\to\R$ satisfying
\begin{equation}\label{eq:continuousstokes}
 \beta(u,v')+\beta(v,u')=\bracket{uv}
\end{equation}
form an affine space with origin
$\beta_0(u,f)=\int_{-1}^{1}uf$ and model space
\[
 \Lambda^2_{\mathrm{cont}}(V/\R)^*.
\]
Explicitly, for every continuous skew form $\Omega$ on $V/\R$,
\begin{equation}\label{eq:continuousbeta}
 \beta_\Omega(u,f)=\int_{-1}^{1}u(x)f(x)\dd x
 +\Omega\bigl([u],\overline d^{-1}f\bigr)
\end{equation}
satisfies \eqref{eq:continuousstokes}, and every compatible continuous pairing has this form.
\end{theorem}

\begin{proof}
If $\beta_1$ and $\beta_0$ are compatible, define
$\Omega(u,v)=(\beta_1-\beta_0)(u,v')$.  Subtraction of the Stokes identities makes $\Omega$
skew, and $v'=0$ for constants, so it descends to $V/\R$.  Conversely,
\eqref{eq:continuousbeta} gives
\[
 \beta_\Omega(u,v')+\beta_\Omega(v,u')
 =\int(uv)'+\Omega([u],[v])+\Omega([v],[u])
 =\bracket{uv}.
\]
Uniqueness follows from surjectivity of $d$.
\end{proof}

Equip $Q^{(1)}=H^1(-1,1)/\R$ with
$\|[u]\|_1=\|u'\|_{L^2}$.  Then $\overline d:Q^{(1)}\to W$ is an isometry.  This derivative
metric is distinct from the $L^2$ quotient metric on $Q_N^{(0)}$ used in
Section~\ref{sec:metric}; the superscripts record that distinction.  Riesz representation
therefore sharpens Theorem~\ref{thm:continuous} as follows.

\begin{proposition}[Operator form and coherent discrete extension]
\label{prop:coherentextension}
Every continuous gauge $\Omega$ has a unique representation
\begin{equation}\label{eq:continuousoperatorform}
 \Omega([u],[v])=\langle u',Sv'\rangle_{L^2},
 \qquad S^*=-S\in\mathcal L(L^2),
\end{equation}
and the compatible pairing is
\begin{equation}\label{eq:continuousbetas}
 \beta_S(u,f)=\int_{-1}^{1}uf\dd x+\langle u',Sf\rangle_{L^2}.
\end{equation}

Let $Q_N^{(1)}=\Pp_N/\R\subset Q^{(1)}$ and let
$\omega_N\in\Lambda^2(Q_N^{(1)})^*$.  There exists a continuous two-form $\Omega$ on
$Q^{(1)}$ with $\Omega|_{Q_N^{(1)}\times Q_N^{(1)}}=\omega_N$ for every $N$ if and only if
\begin{align}
 \omega_{N+1}|_{Q_N^{(1)}\times Q_N^{(1)}}&=\omega_N,
 \label{eq:restrictioncoherence}\\
 \sup_N\|\omega_N\|_{\mathrm{op}}&<\infty,
 \label{eq:uniformcoherence}
\end{align}
where the operator norm is induced by $\|\cdot\|_1$.  The extension, and hence $S$, is
unique.
\end{proposition}

\begin{proof}
The isometry $\overline d$ transports continuous skew forms on $Q^{(1)}$ bijectively to bounded
skew forms on $L^2$; Riesz representation gives the unique skew-adjoint operator $S$ and
\eqref{eq:continuousbetas}.  Necessity of
\eqref{eq:restrictioncoherence}--\eqref{eq:uniformcoherence} is immediate.  Conversely,
restriction compatibility defines a skew form on the increasing union
$\bigcup_NQ_N^{(1)}$, and the uniform bound makes it continuous there.  Polynomials are dense
in $H^1$ in the derivative norm modulo constants, so the form extends uniquely to
$Q^{(1)}\times Q^{(1)}$.
\end{proof}

Thus fixed-degree classification alone does not guarantee a continuum limit: restriction
coherence and uniform operator control are the exact additional requirements.

\subsection{Dirac \texorpdfstring{$B$}{B}-transform}

A two-form $\Omega$ defines
\[
 \tau_\Omega(v,e)=(v,e+\iota_v\Omega)
\]
on $T(V/\R)\oplus T^*(V/\R)$.  Because $\Omega$ is skew, $\tau_\Omega$ preserves the
canonical symmetric power pairing.  Since the form is constant, it is closed.  It therefore
has the algebraic form of a gauge transformation by a closed two-form
\cite{bursztynradko2003}.  This identifies the interior linear transformation.  To conclude
that it transforms a particular open boundary-port Dirac structure, one must additionally
verify its domain, boundary trace relation, and maximal isotropy; those analytical data are
not part of the finite-dimensional affine theorem.  An analytical Hilbert-space treatment of
these domain and boundary-duality questions is given in \cite{brugnoli2023}.

This two-form gauge is related to, but distinct from, the gauge reduction of PHS potentials.
In the latter, a differential form $\rho$ is replaced by $\rho+d\alpha$ and the physical
variable $d\rho$ is unchanged; Stokes--Dirac and simplicial Dirac reductions of this type are
studied in \cite{seslija2012}.  That reduction creates the quotient on which the present
two-form lives; it does not itself vary the mixed pairing.

\subsection{The Le Gorrec--Zwart--Maschke skew freedom}

Le Gorrec, Zwart, and Maschke consider formally skew differential operators of the form
\begin{equation}\label{eq:lgzmoperator}
 J=\sum_{i=0}^{r}P^{(i)}\partial_x^i,
 \qquad
 P^{(i)}=(-1)^{i+1}P^{(i)\mathsf T}.
\end{equation}
In particular, $P^{(0)\mathsf T}=-P^{(0)}$.  Consequently,
\begin{equation}\label{eq:p0cancel}
 \pair{e_1}{P^{(0)}e_2}_{L^2}
 +\pair{e_2}{P^{(0)}e_1}_{L^2}=0,
\end{equation}
and $P^{(0)}$ disappears from the Green boundary form.  Indeed, the boundary matrix $Q$ in
their Theorem~3.1 depends only on $P^{(1)},\ldots,P^{(r)}$
\cite{legorrec2005}.  This is the same power-invisibility principle that produces the present
affine gauge: Stokes fixes the symmetric part of an interior bilinear expression and leaves
its antisymmetric part undetermined.

The precise finite-dimensional relation requires the derivative quotient.  Endow $E_N$ with
the unweighted $L^2$ inner product and regard $dq=q'$ as an element of $E_N$ through the
natural inclusion $F_N\subset E_N$.

\begin{proposition}[Riesz representation of the pairing gauge]\label{prop:relativepzero}
The Riesz map defines a linear isomorphism
\begin{equation}\label{eq:Kspace}
 \Lambda^2(E_N/\R)^*
 \simeq
 \mathfrak k_N
 :=\{K_N\in\operatorname{End}(E_N):K_N^*=-K_N,\ K_N1=0\}.
\end{equation}
In particular,
\begin{equation}\label{eq:Kdimension}
 \dim\mathfrak k_N=\frac{N(N-1)}2.
\end{equation}
More explicitly, every pairing gauge has a unique representation
\begin{equation}\label{eq:omegak}
 \omega(p,q)=\pair{p}{K_Nq}_{L^2},
\end{equation}
and hence
\begin{equation}\label{eq:operatorgauge}
 a_\beta(p,q)
 =\pair{p}{(d+K_N)q}_{L^2}.
\end{equation}
Moreover, there is a unique map $L_N:F_N\to E_N$ such that
\begin{equation}\label{eq:Kfactorization}
 \boxed{K_N=L_Nd.}
\end{equation}
Thus the pairing gauge is exactly a skew-adjoint, zero-power correction to the discrete
interconnection, subject to annihilating the derivative kernel.
\end{proposition}

\begin{proof}
For fixed $q$, the functional $p\mapsto\omega(p,q)$ has a unique $L^2$ Riesz representative,
which defines $K_Nq$ by \eqref{eq:omegak}.  Skew-symmetry of $\omega$ is equivalent to
$K_N^*=-K_N$.  Since $\omega(p,1)=0$ for all $p$, nondegeneracy of $L^2$ gives $K_N1=0$.

Conversely, if $K_N^*=-K_N$ and $K_N1=0$, then \eqref{eq:omegak} is skew and vanishes
whenever either argument is constant; it therefore defines a two-form on $E_N/\R$.  This
proves the isomorphism and \eqref{eq:Kdimension}, since $\dim(E_N/\R)=N$.  Finally, if
$dq_1=dq_2$, then $q_1-q_2$ is constant and
$K_Nq_1=K_Nq_2$.  Hence
\[
 L_N(dq):=K_Nq
\]
is well defined.  Surjectivity of $d:E_N\to F_N$ proves existence and uniqueness of $L_N$,
and gives \eqref{eq:Kfactorization}.
\end{proof}

Proposition~\ref{prop:relativepzero} justifies calling the pairing gauge a
\emph{relative $P^{(0)}$ freedom}: after the $L^2$ identification, it is algebraically an
antisymmetric addition to the interconnection operator, exactly as in
\eqref{eq:lgzmoperator}.  The qualifier ``relative'' is essential.  The rectangular
factorization imposes $\left.K_N\right|_{\Ker d}=0$, a condition absent from a general
$P^{(0)}$ term.

\begin{corollary}[Locality obstruction]\label{cor:p0locality}
Consider the componentwise derivative on
$E_N\otimes\R^s$, whose kernel contains every constant vector field.  Suppose an operator in
the analogue of $\mathfrak k_N$ is the restriction of a pointwise multiplication operator
\begin{equation}\label{eq:localp0}
 (K_Nq)(x)=P^{(0)}(x)q(x),
 \qquad P^{(0)}(x)^{\mathsf T}=-P^{(0)}(x).
\end{equation}
Then $P^{(0)}(x)=0$ almost everywhere.  In particular, every nonzero pairing gauge in the
scalar polynomial setting fails to be a standard local order-zero multiplication term when
transferred to the interconnection operator; its nodal or modal realization is generically
nonlocal.
\end{corollary}

\begin{proof}
For every $c\in\R^s$, the constant function $q(x)=c$ belongs to the derivative kernel.
Therefore $0=K_Nc=P^{(0)}(x)c$ almost everywhere.  Since this holds for every $c$,
$P^{(0)}(x)=0$ almost everywhere.  For $s=1$, the conclusion also follows directly because a
real skew-symmetric $1\times1$ matrix is zero.
\end{proof}

There is a second freedom in the construction of Le Gorrec, Zwart, and Maschke.  If
$R_{\mathrm{ext}}$ factorizes their boundary form, every other admissible factorization is
\begin{equation}\label{eq:boundarygauge}
 R=UR_{\mathrm{ext}},
 \qquad U^{\mathsf T}\Sigma U=\Sigma.
\end{equation}
This is a power-preserving change of boundary-port coordinates.  It leaves the differential
operator and its Green form unchanged and is therefore distinct from the present pairing
gauge, which changes $\beta$, its radical, and generally the roots of $\kappa_\beta$.

We conclude that the three notions require a strict distinction.  The $P^{(0)}$ term expresses
local component-skew freedom in the continuum operator; $K_N$ expresses spatial-mode skew
freedom compatible with the derivative quotient; and $U$ in \eqref{eq:boundarygauge}
expresses only a boundary-coordinate freedom.  They share preservation of power, but they do
not act on the same object.  A nonzero sequence $(K_N)$ can represent a genuine continuum
operator only under additional coherence, locality, and uniform-boundedness hypotheses.

\subsection{Why a change of weight is a different geometry}

Let $w$ be smooth and strictly positive on the closed interval, and define
\[
 g_w(u,v)=\int_{-1}^{1}wuv.
\]
The covariant first-order operator
\begin{equation}\label{eq:weightedderivative}
 \nabla_w=\partial_x+\frac12\frac{w'}w
\end{equation}
satisfies
\begin{equation}\label{eq:weightedstokes}
 g_w(u,\nabla_wv)+g_w(v,\nabla_wu)=\bracket{wuv}.
\end{equation}
Thus changing weight modifies the symmetric metric, the differential operator, and the
boundary form together.  It is conjugate to the unweighted derivative by multiplication with
$\sqrt w$, but it does not belong to the affine class with $d$ and $b$ fixed.

More generally, for a coefficient $\sigma$ define
\[
 \nabla_{\sigma,w}
 =\sigma\partial_x+\frac12\frac{(\sigma w)'}{w}.
\]
Then
\begin{equation}\label{eq:pearsonstokes}
 g_w(u,\nabla_{\sigma,w}v)+g_w(v,\nabla_{\sigma,w}u)
 =\bracket{\sigma wuv}.
\end{equation}
This is the first-order weighted Stokes identity naturally associated with the Pearson pair
$(\sigma,w)$.  Jacobi, Gegenbauer, and Chebyshev weights require the corresponding weighted
domains when $w$ or $\sigma w$ is singular or vanishes at an endpoint; the algebraic identity
then holds on the natural dense test space and the boundary trace is interpreted accordingly.

This distinction resolves the relation with Bochner--Pearson families.  A gauge may be
chosen so that its kernel line coincides with a Chebyshev, Gegenbauer, or symmetric Jacobi
polynomial.  Nevertheless, Theorem~\ref{thm:momentrigidity} prevents the gauged pairing from
being the corresponding weighted moment functional.  To recover the positive orthogonality,
Gaussian quadrature, and Sturm--Liouville operator of the classical family, one must pass to
the weighted Stokes triple \eqref{eq:weightedstokes}.

\section{Relation to existing theory and scope of the claims}

Several neighbouring results are classical or already established.
\begin{itemize}
\item Gaussian nodes are the roots of orthogonal polynomials, and generalized Gaussian
quadrature treats broader Chebyshev systems \cite{szego1975,gautschi2004,karlinstudden1966}.
Positive moment functionals on finite-dimensional function spaces admit finite positive
atomic representations by Tchakaloff's theorem \cite{tchakaloff1957,bayerteichmann2006}.
\item SBP norm matrices carry high-order quadrature information
\cite{hickenzingg2013}; generalized and function-space SBP make quadrature-based existence
and construction explicit \cite{delrey2014,glaubitz2023}, while dual-pairing SBP develops a
different forward--backward operator framework \cite{williamsduru2024}.
\item Recent FSBP work directly constructs optimal nodes through generalized Gaussian and
Gauss--Lobatto quadrature \cite{hale2026,bercik2026}.
\item Dirac structures admit gauge transformations by closed two-forms
\cite{bursztynradko2003}, and Stokes--Dirac systems possess gauge-reduction mechanisms in
continuous and discrete settings \cite{seslija2012}; the Hilbert-space domain and boundary
duality required for maximality are treated analytically in \cite{brugnoli2023}.
\item Weak-form and dual-field finite-element exterior-calculus discretizations already use
power-preserving projection, duality, and representation choices
\cite{kotyczka2018,brugnolirashad2022}.  They provide structured constructions rather than a
classification of all compatible mixed pairings for fixed $(E,F,D,b)$.
\end{itemize}

The present results do not claim a new derivation of Gaussian quadrature from orthogonality,
nor a new general FSBP existence theorem.  The principal claims are instead:
\begin{enumerate}
\item the compatibility and full-rank existence criteria for the abstract length-one Stokes
structure $(E,F,D,b)$, its affine classification, and the rank-two boundary radical theorem;
\item the exact relative-open image, constructive synthesis, fibre dimension, two-component
fibre topology, base topology, global continuous/smooth-section obstruction, minimal atlas,
explicit two-chart synthesis, and radial blow-up theorem for the polynomial projective
kernel map;
\item the hyperbolic barycentre characterization of admissible interior real degree-$N$
root sets;
\item the moment-factorization formulation of locality and the rigidity theorem identifying
$L^2$/Gauss--Legendre as the unique local scalar representative of the polynomial gauge
class;
\item the closed minimum-gauge formula, its exact global rank stratification by kernel degree,
the exact attainment criterion in the full-rank fibre, the rectangular right radical, the
quantitative transverse opening of the lower-degree defect, the
resulting global canonical pairing section on the degree-$N$ kernel stratum, the
canonical nodal section on its simple interior real-root part, its hyperbolic stability near Gauss,
and the collision obstruction to a global estimate;
\item the Riesz identification of the gauge freedom with the skew operators
$K_N^*=-K_N$, $\left.K_N\right|_{\Ker d}=0$, its factorization $K_N=L_Nd$, and the resulting precise
separation from local $P^{(0)}$ terms, boundary-port coordinate gauges, and weighted positive
metrics, together with the exact coherence criterion for extension to $H^1/\R$.
\end{enumerate}
The abstract affine theorem is elementary finite-dimensional linear algebra; generalized
moment representations for Chebyshev systems are classical.  The novelty claimed here lies
in their joint Stokes interpretation and, more specifically, in the exact polynomial kernel
image and fibre theorems, locality rigidity, and the associated gauge/operator distinction.
To the best of the author's knowledge, this combination has not appeared in the SBP,
generalized quadrature, or Stokes--Dirac literature.

\section{Conclusion}

For a surjective length-one discrete Stokes structure, compatibility is equivalent to the
vanishing of the boundary form on $\Ker D\times\Ker D$.  The compatible pairings then form
an affine space directed by $\Lambda^2(E/\Ker D)^*$; full rank is generic whenever it is
allowed by the parity--boundary criterion.  Stokes therefore fixes the symmetric boundary
power but leaves a genuine two-form freedom on the derivative quotient.

For $\Pp_N\xrightarrow{d/dx}\Pp_{N-1}$, the radical of a full-rank pairing is an intrinsic,
gauge-dependent polynomial line.  The kernel map has an exact relative semialgebraic image,
constructive fibres of dimension $(N-1)(N-2)/2$, and two Pfaffian components for $N\geq3$.
On the simple interior real-root stratum, admissibility becomes
\[
 \sum_i\operatorname{artanh}(x_i)=0,
\]
so admissible node sets form a codimension-one family whose equal-weight hyperbolic
barycentre is the origin.

These intrinsic roots are not automatically quadrature nodes.  They provide exact nodal
coordinates with a full coupling matrix.  Scalar locality is a separate moment condition;
because differentiated products span $\Pp_{2N-1}$, it collapses the affine class to the
unweighted $L^2$ pairing and uniquely selects the Gauss--Legendre nodes and positive Gauss
weights.  Thus Gauss--Legendre is not imposed by Stokes alone, but is rigid within the fixed
unweighted Stokes geometry once local scalar power evaluation is required.

After the $L^2$ quotient metric is chosen, the minimum-norm solution of a prescribed-kernel
constraint is represented by the explicit rank-two operator \eqref{eq:Kcanonical}.  The
associated mixed pairing has rank exactly $N$ on the degree-$N$ stratum and $N-2$ on every
nonempty lower-degree stratum, where its full left and right radicals are explicit.  The
defect opens transversely under a Legendre-direction perturbation.  On the degree-$N$
stratum this gives a global canonical pairing section, whose nodal restriction is locally
Lipschitz near Gauss but cannot satisfy a uniform global estimate without a node-separation
or conditioning factor.

The normalized kernel image is $\R^{N-1}$ in odd degree and
$\R^{N-1}\setminus\{0\}$ in even degree.  Although every admissible kernel has an individual
full-rank realization, a global smooth realization over the entire image exists only for odd
$N$ and for $N=2,4,8$.  In even degree this is precisely the almost-complex obstruction on
$S^{N-2}$.  The case $N=2$ is trivial, while the exceptional sections for $N=4$ and $N=8$
come from the three- and seven-dimensional vector cross products.  At all other even orders,
one global chart is impossible but the explicit Rodrigues construction gives an optimal
two-chart atlas.  This is a topological obstruction to choosing the gauge, not an obstruction
to the existence of the fibres.  Independently of that topology, every even-order synthesis
is analytically unbounded near the deleted constant class: the radial equation forces
$\|r_z\|\geq2/\|z\|$ and hence $\|M_z\|_{\mathrm F}\gtrsim\|z\|^{-1}$.  The obstruction is
therefore present even in the exceptional orders $N=2,4,8$.

Finally, the discrete gauge is a quotient-compatible, generally nonlocal skew operator that
factors through differentiation.  Coherent uniformly bounded families extend to the
continuous $H^1/\R$ model.  This separates the pairing gauge from boundary-port coordinate
changes, local $P^{(0)}$ terms, and weighted Hamiltonian geometries.  Classifying which
coherent nonlocal gauges survive the continuous limit is the natural next question.

\appendix

\section{Basis covariance}

Let $T_E$ and $T_F$ be nonsingular changes of basis in $E_N$ and $F_N$.  The mixed and
derivative matrices transform as
\[
 M\mapsto T_E^{\mathsf T}MT_F,
 \qquad
 D\mapsto T_F^{-1}DT_E.
\]
Therefore
\[
 MD\mapsto T_E^{\mathsf T}(MD)T_E,
\]
and the gauge two-form transforms by congruence.  The coefficient vector of
$\kappa_\beta$ changes contragrediently, but the polynomial and its zero set do not.  The
affine classification, endpoint condition, and kernel map are basis-independent.

\section{Constructive synthesis for a prescribed kernel}

Given a nonconstant $\kappa\in\Pp_N$ satisfying
$\kappa(1)=(-1)^N\kappa(-1)\neq0$, the proof of
Theorem~\ref{thm:kernelimage} gives an explicit construction.
For $N=1$, the only projective kernel is $[x]$ and the unique compatible pairing is the
zero-gauge $L^2$ representative.  Assume below that $N\geq2$.
\begin{enumerate}
\item Decompose $\kappa=\alpha+\beta x+z$ with
$z\in(1-x^2)\Pp_{N-2}$.
\item If $N$ is even, choose a skew matrix $C$ on $Z_N$ with
$\Ker C=\spanof\{z\}$ and choose $r$ satisfying $z^{\mathsf T}r=2\alpha$.
\item If $N$ is odd, choose any invertible skew matrix $C$ and set
$r=\beta^{-1}Cz$.
\item Form $A_{r,C}$ from \eqref{eq:Anormal}.  Because
$A_{r,C}(\cdot,1)=0$, the form factors through the derivative in its second slot; surjectivity
of $d:E_N\to F_N$ then gives a unique mixed matrix $M$ such that $MD=A_{r,C}$.
\end{enumerate}
The resulting $M$ is full rank, Stokes-compatible, and has left radical
$\spanof\{\kappa\}$.  This construction turns the kernel image theorem into an exact
synthesis procedure rather than a dimension count.

\begin{thebibliography}{99}
\footnotesize

\bibitem{baez2002}
J.~C.~Baez,
\newblock The octonions,
\newblock \emph{Bulletin of the American Mathematical Society}, 39 (2002), pp.~145--205.
\newblock \url{https://doi.org/10.1090/S0273-0979-01-00934-X}.

\bibitem{bayerteichmann2006}
C.~Bayer and J.~Teichmann,
\newblock The proof of Tchakaloff's theorem,
\newblock \emph{Proceedings of the American Mathematical Society}, 134 (2006),
pp.~3035--3040.
\newblock \url{https://doi.org/10.1090/S0002-9939-06-08249-9}.

\bibitem{bercik2026}
A.~Bercik, L.~Patrascu, and D.~W.~Zingg,
\newblock Construction and optimization of summation-by-parts operators for general function spaces using an improved generalized Gaussian quadrature algorithm,
\newblock \emph{arXiv preprint} arXiv:2607.08934, 2026.
\newblock \url{https://doi.org/10.48550/arXiv.2607.08934}.

\bibitem{borelserre1953}
A.~Borel and J.-P.~Serre,
\newblock Groupes de Lie et puissances r\'eduites de Steenrod,
\newblock \emph{American Journal of Mathematics}, 75 (1953), pp.~409--448.
\newblock \url{https://doi.org/10.2307/2372495}.

\bibitem{browngray1967}
R.~B.~Brown and A.~Gray,
\newblock Vector cross products,
\newblock \emph{Commentarii Mathematici Helvetici}, 42 (1967), pp.~222--236.
\newblock \url{https://doi.org/10.1007/BF02564418}.

\bibitem{brugnolirashad2022}
A.~Brugnoli, R.~Rashad, and S.~Stramigioli,
\newblock Dual field structure-preserving discretization of port-Hamiltonian systems using
finite element exterior calculus,
\newblock \emph{Journal of Computational Physics}, 471 (2022), article 111601.
\newblock \url{https://doi.org/10.1016/j.jcp.2022.111601}.

\bibitem{brugnoli2023}
A.~Brugnoli, G.~Haine, and D.~Matignon,
\newblock Stokes--Dirac structures for distributed parameter port-Hamiltonian systems:
an analytical viewpoint,
\newblock \emph{Communications in Analysis and Mechanics}, 15 (2023), pp.~362--387.
\newblock \url{https://doi.org/10.3934/cam.2023018}.

\bibitem{bursztynradko2003}
H.~Bursztyn and O.~Radko,
\newblock Gauge equivalence of Dirac structures and symplectic groupoids,
\newblock \emph{Annales de l'Institut Fourier}, 53 (2003), pp.~309--337.
\newblock \url{https://doi.org/10.5802/aif.1945}.

\bibitem{cheng1999}
H.~Cheng, V.~Rokhlin, and N.~Yarvin,
\newblock Nonlinear optimization, quadrature, and interpolation,
\newblock \emph{SIAM Journal on Optimization}, 9 (1999), pp.~901--923.
\newblock \url{https://doi.org/10.1137/S1052623498349796}.

\bibitem{delrey2014}
D.~C.~Del Rey Fern\'andez, P.~D.~Boom, and D.~W.~Zingg,
\newblock A generalized framework for nodal first derivative summation-by-parts operators,
\newblock \emph{Journal of Computational Physics}, 266 (2014), pp.~214--239.
\newblock \url{https://doi.org/10.1016/j.jcp.2014.01.038}.

\bibitem{gautschi2004}
W.~Gautschi,
\newblock \emph{Orthogonal Polynomials: Computation and Approximation},
\newblock Oxford University Press, Oxford, 2004.
\newblock \url{https://doi.org/10.1093/oso/9780198506720.001.0001}.

\bibitem{glaubitz2023}
J.~Glaubitz, J.~Nordstr\"om, and P.~\"Offner,
\newblock Summation-by-parts operators for general function spaces,
\newblock \emph{SIAM Journal on Numerical Analysis}, 61 (2023), pp.~733--754.
\newblock \url{https://doi.org/10.1137/22M1470141}.

\bibitem{gustafson1970}
S.-\AA.~Gustafson,
\newblock On the computational solution of a class of generalized moment problems,
\newblock \emph{SIAM Journal on Numerical Analysis}, 7 (1970), pp.~343--357.
\newblock \url{https://doi.org/10.1137/0707026}.

\bibitem{hale2026}
N.~Hale, C.~Harley, P.~Nchupang, and J.~Nordstr\"om,
\newblock Summation-by-parts operators for general function spaces: optimal nodes,
\newblock \emph{arXiv preprint} arXiv:2604.23306, 2026.
\newblock \url{https://doi.org/10.48550/arXiv.2604.23306}.

\bibitem{hickenzingg2013}
J.~E.~Hicken and D.~W.~Zingg,
\newblock Summation-by-parts operators and high-order quadrature,
\newblock \emph{Journal of Computational and Applied Mathematics}, 237 (2013), pp.~111--125.
\newblock \url{https://doi.org/10.1016/j.cam.2012.07.015}.

\bibitem{hirsch1976}
M.~W.~Hirsch,
\newblock \emph{Differential Topology},
\newblock Graduate Texts in Mathematics 33, Springer, New York, 1976.
\newblock \url{https://doi.org/10.1007/978-1-4684-9449-5}.

\bibitem{karlinstudden1966}
S.~Karlin and W.~J.~Studden,
\newblock \emph{Tchebycheff Systems: With Applications in Analysis and Statistics},
\newblock Interscience, New York, 1966.

\bibitem{konstantisparton2017}
P.~Konstantis and M.~Parton,
\newblock Almost complex structures on spheres,
\newblock \emph{Differential Geometry and its Applications}, 57 (2018), pp.~10--22.
\newblock \url{https://doi.org/10.1016/j.difgeo.2017.10.010}.

\bibitem{kotyczka2018}
P.~Kotyczka, B.~Maschke, and L.~Lef\`evre,
\newblock Weak form of Stokes--Dirac structures and geometric discretization of
port-Hamiltonian systems,
\newblock \emph{Journal of Computational Physics}, 361 (2018), pp.~442--476.
\newblock \url{https://doi.org/10.1016/j.jcp.2018.02.006}.

\bibitem{legorrec2005}
Y.~Le Gorrec, H.~Zwart, and B.~Maschke,
\newblock Dirac structures and boundary control systems associated with skew-symmetric differential operators,
\newblock \emph{SIAM Journal on Control and Optimization}, 44 (2005), pp.~1864--1892.
\newblock \url{https://doi.org/10.1137/040611677}.

\bibitem{micchellipinkus1977}
C.~A.~Micchelli and A.~Pinkus,
\newblock Moment theory for weak Chebyshev systems with applications to monosplines,
quadrature formulae and best one-sided $L^1$-approximation by spline functions with fixed
knots,
\newblock \emph{SIAM Journal on Mathematical Analysis}, 8 (1977), pp.~206--230.
\newblock \url{https://doi.org/10.1137/0508015}.

\bibitem{moulla2012}
R.~Moulla, L.~Lef\`evre, and B.~Maschke,
\newblock Pseudo-spectral methods for the spatial symplectic reduction of open systems of conservation laws,
\newblock \emph{Journal of Computational Physics}, 231 (2012), pp.~1272--1292.
\newblock \url{https://doi.org/10.1016/j.jcp.2011.10.008}.

\bibitem{seslija2012}
M.~Seslija, A.~J.~van der Schaft, and J.~M.~A.~Scherpen,
\newblock Reduction of Stokes--Dirac structures and gauge symmetry in port-Hamiltonian systems,
\newblock \emph{IFAC Proceedings Volumes}, 45 (2012), pp.~114--119.
\newblock \url{https://doi.org/10.3182/20120829-3-IT-4022.00030}.

\bibitem{szego1975}
G.~Szeg\H{o},
\newblock \emph{Orthogonal Polynomials}, 4th ed.,
\newblock American Mathematical Society, Providence, RI, 1975.

\bibitem{tchakaloff1957}
V.~Tchakaloff,
\newblock Formules de cubatures m\'ecaniques \`a coefficients non n\'egatifs,
\newblock \emph{Bulletin des Sciences Math\'ematiques}, 81 (1957), pp.~123--134.

\bibitem{vanderschaftmaschke2002}
A.~J.~van der Schaft and B.~M.~Maschke,
\newblock Hamiltonian formulation of distributed-parameter systems with boundary energy flow,
\newblock \emph{Journal of Geometry and Physics}, 42 (2002), pp.~166--194.
\newblock \url{https://doi.org/10.1016/S0393-0440(01)00083-3}.

\bibitem{williamsduru2024}
C.~Williams and K.~Duru,
\newblock Full-spectrum dispersion relation preserving summation-by-parts operators,
\newblock \emph{SIAM Journal on Numerical Analysis}, 62 (2024), pp.~1565--1588.
\newblock \url{https://doi.org/10.1137/23M1586471}.

\end{thebibliography}

\end{document}
